\documentclass[11pt,a4paper]{article}
\usepackage[T1]{fontenc}
\usepackage[utf8]{inputenc}
\usepackage[french]{babel}
\frenchsetup{AutoSpacePunctuation=false}
\addto\captionsfrench{\def\tablename{Tableau}}
\usepackage{lmodern}
\usepackage[margin=2.6cm]{geometry}
\usepackage{amsmath,amssymb,amsthm}
\usepackage{booktabs}
\usepackage{array}
\usepackage[hidelinks]{hyperref}

\newtheorem{theorem}{Théorème}[section]
\newtheorem{lemma}[theorem]{Lemme}
\newtheorem{proposition}[theorem]{Proposition}
\newtheorem{corollary}[theorem]{Corollaire}
\theoremstyle{definition}
\newtheorem{definition}[theorem]{Définition}
\theoremstyle{remark}
\newtheorem{remark}[theorem]{Remarque}

\newcommand{\Q}{\mathbb{Q}}
\newcommand{\E}{\mathbb{E}}
\newcommand{\PP}{\mathbb{P}}
\newcommand{\cR}{\mathcal{R}}
\newcommand{\cI}{\mathcal{I}}
\newcommand{\sep}{\operatorname{sep}}
\newcommand{\sepw}{\operatorname{sep}_{\theta}}
\newcommand{\dfc}{\operatorname{def}}
\newcommand{\tw}{\operatorname{tw}}
\newcommand{\rank}{\operatorname{rank}}
\newcommand{\Gr}{\Gamma}
\newcommand{\nocc}{n_{\mathrm{occ}}}
\newcommand{\Nbar}{\overline{N}}
\newcommand{\DL}{DL}

% Citations de Dias--Lima : reproduites en anglais (césure anglaise), puis traduites.
\makeatletter
\newcommand{\englishhyph}{\@ifundefined{l@english}{}{\language\l@english}}
\makeatother
\newcommand{\eng}[1]{{\englishhyph #1}}
\newenvironment{dlquote}{\begin{quote}\englishhyph}{\end{quote}}
\newenvironment{dltrad}{\begin{quote}\itshape\textup{Traduction~:}\ \ignorespaces}{\end{quote}}

\title{Les lignes de mesure absorbent au plus\\ un facteur \texorpdfstring{$\sqrt{s}$}{sqrt(s)} de la coupe :\\
une borne par déterminant pour les schémas\\ d'interpolation par produit de
clauses\thanks{Version française de l'article original en anglais ; en cas de divergence, la
version anglaise fait foi.}}
\author{Lucas Requena\\[2pt]
{\small Étudiant en informatique}\\
{\small \texttt{l.requena@outlook.fr}}}
\date{Novembre 2026}

\begin{document}
\maketitle

\begin{abstract}
Dias et Lima (arXiv:2611.00517) abstraient la forme produit de clauses de l'algorithme
d'interpolation figurant dans la réfutation, revendiquée par OpenAI, de l'hypothèse du temps
exponentiel (ETH) sous la forme d'un \emph{schéma} : en chaque point d'évaluation, les
variables sont réparties en groupes de taille au plus $s$, et une clause ne peut voir les
autres groupes qu'à travers des \emph{lignes de mesure} linéaires. Ils démontrent un coût de
$2^{\Omega(\sqrt n)}$ sur les formules à expansion et laissent ouverte, dans leur
Section~9, la question de savoir si les mesures peuvent \eng{``absorb the
cut''}\footnote{Traduction : \og absorber la coupe\fg{}.}. Nous montrons que les mesures en
absorbent au plus un facteur $\sqrt{s}$ : en un certain point, le nombre de lignes plus
l'ordre de Taylor dépasse une constante fois $\sep(F)/\sqrt{s}$, où $\sep(F)$ est le nombre
de séparation équilibrée du graphe primal. Pour les déformations qui ne dépendent que de la
variable, le facteur $\sqrt s$ disparaît. La démonstration utilise un déterminant
$n\times n$ construit à partir d'une occurrence par variable, ainsi qu'une borne
d'anticoncentration. En conséquence, sur les formules à occurrences bornées à séparateurs
linéaires et sur les formules 3-CNF aléatoires de densité supérieure à $3$, tout schéma
appliqué à la formule elle-même coûte $2^{\Omega(n^{2/3})}$ dans le modèle de coût de type
produit. Nous observons aussi que le modèle de coût de type PD (programmation dynamique) est
dégénéré : un schéma trivial a une charge de $2^{\lceil\sqrt n\rceil}$ sur toute formule. La
partie déterministe de la démonstration est vérifiée par machine en Lean~4. Les résultats
ont été obtenus avec des outils d'IA, n'ont pas encore été relus par des experts humains, et
ne disent rien des voies plus fines de la prépublication, ni d'ETH, ni de P contre NP.
\end{abstract}

\section{Introduction}\label{sec:intro}

L'hypothèse du temps exponentiel (ETH) d'Impagliazzo et Paturi~\cite{IP01} affirme que 3-SAT
sur $n$ variables nécessite un temps $2^{\delta n}$ pour un certain $\delta>0$. En octobre
2026, une prépublication d'OpenAI~\cite{OpenAI} a revendiqué un algorithme déterministe qui
compte les affectations satisfaisantes d'une formule $k$-CNF en temps
$2^{O_k(n\log\log n/\sqrt{\log n})}$, ce qui réfuterait ETH. Un exposé en a été rédigé par
O'Donnell~\cite{ODonnell}, et une seconde prépublication, due à Claude
d'Anthropic~\cite{Claude}, a annoncé un algorithme en temps $2^{o(n)}$ pour 3-SAT. Nous
n'évaluons la correction d'aucune de ces deux prépublications, et nous n'étudions pas la
seconde.

Dias et Lima~\cite{DL} ont analysé la prépublication d'OpenAI. Entre autres choses, ils
isolent la \emph{forme produit de clauses} de son algorithme d'interpolation et l'abstraient
en un \emph{schéma d'interpolation par produit de clauses} (leur Définition~5.1, reproduite
à la Section~\ref{sec:setting}). Dans un modèle de coût qui facture ce que paie l'évaluation
de la prépublication, ils démontrent que tout schéma de ce type coûte
$2^{\Omega(\sqrt{\tw})}$ sur les formules de largeur arborescente primale $\tw$, donc
$2^{\Omega(\sqrt n)}$ sur les formules à expansion. Le design par produit de clauses de la
prépublication elle-même coûte $2^{O(n(\log\log n)^2/\sqrt{\log n})}$ sur les formules à
occurrences bornées ; il reste donc un écart. Dans leur Section~9, Dias et Lima expliquent
que leur technique ne peut pas aller au-delà de $2^{\tilde O(\sqrt n)}$, et ils énoncent ce
qu'il faut pour réduire l'écart par le bas : \eng{``one must show that measurements cannot
absorb the cut''}\footnote{Traduction : \og il faut montrer que les mesures ne peuvent pas
absorber la coupe\fg{}.}. Ils font une observation par paires pour les groupes ayant une
seule ligne de mesure (un mineur $2\times2$) et écrivent que \eng{``turning this into a
bound for groups with several measurement rows is open''}\footnote{Traduction : \og
transformer cela en une borne pour des groupes ayant plusieurs lignes de mesure est un
problème ouvert\fg{}.}.

\paragraph{Notre résultat.} Nous donnons une telle borne. Soit $F$ une formule CNF sur $n$
variables dans laquelle chaque variable apparaît au plus $b$ fois, et soit $\sep(F)$ la
taille minimale d'un ensemble de variables dont la suppression laisse des composantes d'au
plus $2n/3$ sommets dans le graphe primal. Avec les notations de la Définition~5.1
de~\cite{DL}, soit $s_t$ la plus grande taille de groupe au point $t$, soit $\rho_t$ le
nombre total de lignes de mesure indépendantes, et soit $q_t$ l'ordre de Taylor. Le
Théorème~\ref{thm:general} affirme que tout schéma pour $F$ possède un point $t$ tel que
\begin{equation}\label{eq:intro}
\rho_t+q_t+1\;>\;\min\Bigl(c_b\,\frac{\sep(F)}{\sqrt{s_t}},\ \frac{2n}{3}-s_t\Bigr),
\qquad c_b=\frac{(1-e^{-2})(b-1)}{4\sqrt2\,b^2},\quad c_3=\frac{1-e^{-2}}{18\sqrt2}=0.03396\ldots
\end{equation}
Si la déformation d'une occurrence ne dépend que de sa variable et non de sa clause, le
Théorème~\ref{thm:uniform} donne le même énoncé avec $\sep(F)$ à la place de
$c_b\sep(F)/\sqrt{s_t}$. Ainsi, lorsque $\sep(F)=\Omega(n)$, un schéma a besoin d'environ
une ligne de mesure (ou une unité d'ordre de Taylor) pour $\sqrt{s_t}$ variables de coupe en
général, et d'une par variable de coupe dans le cas uniforme par variable. Nous ne savons pas
si le facteur $\sqrt{s_t}$ peut réellement être gagné par un design dépendant de la clause ;
cette question est ouverte.

\paragraph{Contributions.}
\begin{enumerate}
\item \textbf{Un lemme du déterminant témoin} (Lemme~\ref{lem:det}). On choisit une
 occurrence de chaque variable et l'on forme la matrice polynomiale $n\times n$ dont les
 lignes sont les vecteurs de déformation choisis. Son déterminant est un polynôme non nul de
 degré au plus $nR$, et son ordre d'annulation en un point d'évaluation est au moins égal à
 un \emph{déficit} combinatoire du design en ce point, moins le nombre de lignes de mesure.
 Le mineur $2\times2$ de \cite[Section~9]{DL} est un mineur $2\times 2$ de cette matrice ;
 le déterminant complet traite un nombre quelconque de lignes par groupe.
\item \textbf{Les déformations uniformes par variable ne gagnent rien aux mesures}
 (Théorème~\ref{thm:uniform}) : un certain point vérifie
 $\rho_t+q_t+1>\min(\sep(F),2n/3-s_t)$.
\item \textbf{Les déformations générales gagnent au plus un facteur $\sqrt{s}$}
 (Théorème~\ref{thm:general}), en moyennant le lemme du déterminant sur des choix aléatoires
 d'occurrences et en minorant l'espérance du déficit par une inégalité d'anticoncentration à
 constante explicite (Lemmes~\ref{lem:anti} et~\ref{lem:expdef}).
\item \textbf{Conséquences} (Corollaire~\ref{cor:main}) pour les formules à occurrences
 bornées avec $\sep(F)\ge cn$ et pour les formules 3-CNF aléatoires à $m=\Delta n$ clauses,
 $\Delta>3$ : le design décrit dans \cite[Section~9]{DL} (taille des groupes, ordre de
 Taylor et nombre de lignes tous en $\tilde O(\sqrt n)$) n'existe pas, et dans le modèle de
 coût de type produit de \cite[Section~5.2]{DL} tout schéma appliqué à la formule elle-même
 coûte $2^{\Omega(n^{2/3})}$ ; les schémas uniformes par variable coûtent $2^{\Omega(n)}$.
\item \textbf{Le modèle de coût de type PD est dégénéré} (Proposition~\ref{prop:dp}). Dans le
 second modèle de coût de \cite[Section~5.2]{DL}, un schéma trivial (déformation constante,
 groupes entièrement mesurés de taille $\lceil\sqrt n\rceil$) a une charge de
 $2^{\lceil\sqrt n\rceil}$ sur toute formule. La borne $2^{\Omega(\sqrt n)}$ de~\cite{DL}
 est donc optimale dans ce modèle pour une raison triviale, aucune technique ne peut y
 démontrer davantage, et la question de la Section~9 n'a de sens que dans un modèle qui
 facture les lignes, comme le modèle de type produit.
\item \textbf{Vérification par machine} (Section~\ref{sec:lean}). Le lemme du déterminant,
 l'étape de comptage de Hermite, l'étape de passage à la moyenne, le lemme du séparateur, le
 Théorème~\ref{thm:uniform}, et la déduction du Théorème~\ref{thm:general} à partir du lemme
 sur l'espérance du déficit sont formalisés en Lean~4 avec Mathlib. Le lemme sur l'espérance
 du déficit lui-même, le lemme sur les 3-CNF aléatoires et le corollaire ne sont démontrés
 que sur papier.
\item \textbf{Vérifications par le calcul} (Section~\ref{sec:checks}) sur de petites instances,
 qui peuvent détecter des erreurs mais ne démontrent rien d'asymptotique.
\end{enumerate}

\paragraph{Ce que cet article ne montre pas.}
\begin{itemize}
\item Les énoncés de coût sont limités au modèle de coût \emph{de type produit}
 de~\cite[Section~5.2]{DL}. Ce sont des bornes inférieures sur les charges de ce modèle, et
 ce ne sont des énoncés sur le temps d'exécution que dans la mesure où ces charges sont
 payées. Dans le modèle de type PD, rien ne peut être démontré au-delà de
 $2^{\lceil\sqrt n\rceil}$ (Proposition~\ref{prop:dp}).
\item Les résultats concernent un schéma appliqué à $F$ \emph{elle-même}. La prépublication
 applique son schéma à une formule à occurrences bornées $F^\star$ obtenue à partir de $F$
 par des copies de variables, et aux feuilles d'une sparsification. Le Théorème~8.3
 de~\cite{DL} couvre ces cas ; notre Corollaire~\ref{cor:main} ne les couvre pas, car nous
 n'avons pas démontré $\sep(F^\star)=\Omega(n^\star)$ pour ces formules.
\item Les résultats ne disent rien de la forme de parité de la prépublication ni de sa
 contraction de tenseurs, qui sortent du cadre de la Définition~5.1 (voir
 \cite[Section~5.1]{DL}), rien de l'algorithme de~\cite{Claude}, et rien de la vérité d'ETH.
\item Les résultats n'ont aucune incidence sur P contre NP. Ce sont des barrières pour une
 abstraction d'une voie algorithmique, et non pour un problème.
\item Nous ne donnons aucune borne supérieure. La question de savoir si un design atteint
 $\rho_t\approx n/\sqrt{s_t}$, et donc si $n^{2/3}$ est le bon exposant dans le modèle de
 type produit, est ouverte.
\item Aucun expert humain n'a encore relu ce travail (Section~\ref{sec:ai}).
\end{itemize}

\paragraph{L'idée en termes simples.}
En un point d'évaluation, les variables sont réparties en groupes. Une clause affectée à un
groupe ne peut lire les variables d'un autre groupe qu'à travers les lignes de mesure de ce
groupe. Choisissons, pour chaque variable $v$, une clause qui la contient, et regardons la
forme linéaire que cette clause lit à la place de $x_v$. En $T=0$, cette forme est $x_v$
elle-même, donc les $n$ formes sont indépendantes et leur déterminant vaut $1$. En un point
d'évaluation $t$, considérons un groupe $S$ à $a$ variables. Les formes attachées à des
clauses d'autres groupes ne peuvent apporter, sur les coordonnées de $S$, que des vecteurs de
l'espace des lignes de $S$. Ainsi, si moins de $a$ des clauses choisies sont affectées à
$S$, et si $S$ a peu de lignes, les $n$ formes deviennent dépendantes en $t$ : le
déterminant s'y annule, à un ordre qui compte les dimensions manquantes. Le déterminant est
de degré au plus $MR$, où $M$ est le nombre d'occurrences, donc la somme de ses ordres
d'annulation aux points d'évaluation est au plus $MR$. D'autre part, une règle de
reconstruction qui est exacte pour tous les polynômes de degré $MR$ doit utiliser plus de
$MR$ données de Hermite au total. Par conséquent, en un certain point, les dimensions
manquantes, c'est-à-dire le \emph{déficit}, doivent être payées par des lignes ou par de
l'ordre de Taylor.

Quelle est la taille du déficit ? Si la déformation ne dépend que de la variable, toute
clause contenant $v$ lit la même forme, et une variable dont les clauses sont affectées à
deux groupes est contrainte partout ; le déficit est alors le nombre de ces variables
\emph{scindées}, et celles-ci forment un séparateur du graphe primal. Si la déformation
dépend de la clause, un design peut espérer que les clauses choisies remplissent exactement
chaque groupe. Mais le choix des clauses nous appartient : si on les choisit au hasard, le
nombre de clauses choisies affectées à un groupe fluctue d'environ la racine carrée de sa
variance, et ces fluctuations s'additionnent pour donner au moins (nombre de variables
scindées)$/\sqrt{s}$. C'est de là que vient le $\sqrt s$. Il s'agit d'une perte due à la
technique de démonstration, et nous ne savons pas si elle reflète une possibilité réelle.

\paragraph{Organisation.} La Section~\ref{sec:related} discute les travaux connexes et
compare nos bornes à celles de~\cite{DL}. La Section~\ref{sec:setting} reproduit les
définitions de~\cite{DL} et fixe les notations. La Section~\ref{sec:results} énonce les
résultats et la Section~\ref{sec:proofs} les démontre. La Section~\ref{sec:dp} traite le
modèle de type PD. Les Sections~\ref{sec:lean} et~\ref{sec:checks} décrivent la
formalisation en Lean et les vérifications par le calcul. La Section~\ref{sec:ai} décrit la
manière dont le travail a été réalisé, et la Section~\ref{sec:concl} conclut par des
problèmes ouverts. L'Annexe~\ref{app:random} reproduit, par souci de complétude, les
démonstrations de deux faits sur les 3-CNF aléatoires tirés de~\cite{DL}.

\section{Travaux connexes}\label{sec:related}

\paragraph{ETH et les prépublications de 2026.} ETH a été formulée par Impagliazzo et
Paturi~\cite{IP01}, et le lemme de sparsification d'Impagliazzo, Paturi et Zane~\cite{IPZ01}
la transfère à de nombreux problèmes. La prépublication d'OpenAI~\cite{OpenAI} revendique un
comptage déterministe sous-exponentiel pour les $k$-CNF. D'après~\cite{DL}, sa borne
principale est démontrée par une contraction de réseaux de tenseurs booléens, et le corps du
texte développe un algorithme d'interpolation sous deux formes, une forme produit de clauses
et une forme de parité. O'Donnell~\cite{ODonnell} a rédigé un exposé de la prépublication.
Une prépublication distincte, due à Claude d'Anthropic~\cite{Claude}, annonce un algorithme
en temps $2^{o(n)}$ pour 3-SAT. Nos résultats ne concernent ni la contraction de tenseurs,
ni la forme de parité, ni l'algorithme de~\cite{Claude}.

\paragraph{Dias et Lima.} Notre point de départ est~\cite{DL}. Toutes nos définitions de
schémas et de modèles de coût sont les leurs. Nous utilisons les résultats suivants
de~\cite{DL}, ou nous nous y comparons.
\begin{itemize}
\item \emph{Lemme~6.2 (produit qui s'annule).} Le produit $\prod_o c_o(T)$ des coefficients
 propres est de degré au plus $MR$ et vaut $1$ en $T=0$, donc en un certain point $t$ au
 plus $q_t$ occurrences perdent leur propre variable. Notre Lemme~\ref{lem:hermite} utilise
 le même argument, appliqué au polynôme $\prod_{t}(T-t)^{q_t+1}$ ; l'argument est le leur.
 Notre Lemme~\ref{lem:det} remplace le produit des $M$ coefficients diagonaux $c_o$ par le
 déterminant d'une matrice $n\times n$, de degré au plus $nR$.
\item \emph{Théorème~7.1(ii).} Au point de leur Lemme~6.2, la suppression d'au plus
 $q_t+b\,\mu_ts_t$ arêtes du graphe d'incidence laisse chaque composante à l'intérieur d'un
 seul groupe, où $\mu_t$ est le nombre de groupes mesurés. Si $s_t\le 2n/3$, les variables
 incidentes aux arêtes supprimées forment un $2/3$-séparateur du graphe
 primal\footnote{Si $u,v$ sont adjacentes dans le graphe primal et qu'aucune des deux n'est
 incidente à une arête supprimée, elles partagent une clause $e$ et les deux arêtes
 $(e,u),(e,v)$ subsistent, donc $u,v$ sont dans la même composante après la suppression,
 donc dans le même groupe.}, ce qui donne
 \[
 q_t+b\,\mu_t s_t\ \ge\ \sep(F).
 \]
 Cette borne facture les \emph{groupes mesurés} : un groupe mesuré de taille $s_t$ peut
 absorber la totalité de ses quelque $bs_t$ occurrences étrangères, quel que soit son nombre
 de lignes. Comme $\rho_t\ge\mu_t$, elle implique $\rho_t\ge(\sep(F)-q_t)/(bs_t)$. Notre
 Théorème~\ref{thm:general} facture les \emph{lignes} et donne
 $\rho_t>c_b\sep(F)/\sqrt{s_t}-q_t-1$ lorsque $c_b\sep(F)/\sqrt{s_t}\le 2n/3-s_t$. En bref :
 $\mu_ts_t\gtrsim n$ dans~\cite{DL}, contre $\rho_t\sqrt{s_t}\gtrsim n$ ici, sur les formules
 avec $\sep(F)=\Omega(n)$ et en des points $t$ éventuellement différents.
\item \emph{Théorème~8.1 (barrière de largeur arborescente) et Lemme~8.2.} Dans le modèle de
 coût, tout schéma de type PD, donc tout schéma de type produit, sur $F$ coûte au moins
 $2^{\min(\sqrt{(\tw+1)/2}-1,\,(\tw+1)/2,\,\lfloor n/2\rfloor)}$, où $\tw$ est la largeur
 arborescente primale, et une 3-CNF aléatoire avec $\Delta>3$ a une largeur arborescente
 primale linéaire avec forte probabilité. Un graphe de largeur arborescente $\tw$ possède un
 séparateur d'au plus $\tw+1$ sommets laissant des composantes d'au plus $n/2$
 sommets~\cite{RS86,Bod93}, donc $\sep(F)\le\tw+1$, et le Théorème~8.1 de~\cite{DL}
 s'applique aux familles de notre Corollaire~\ref{cor:main}. Il donne $2^{\Omega(\sqrt n)}$
 pour un schéma appliqué à $F$ elle-même. C'est la comparaison équitable pour notre
 $2^{\Omega(n^{2/3})}$, qui ne vaut que dans le modèle de type produit.
 Notre Lemme~\ref{lem:random} adapte la démonstration de leur Lemme~8.2 de la largeur
 arborescente aux $2/3$-séparateurs et utilise les deux faits de leur Annexe~A.
\item \emph{Théorème~8.3 (3-CNF aléatoires).} Avec forte probabilité, le schéma par produit
 de clauses coûte $2^{\Omega_\Delta(\sqrt n)}$ lorsqu'il est appliqué (a) à la formule à
 occurrences bornées $F^\star$ de la prépublication et (b) aux feuilles de sa
 sparsification. Nos résultats n'améliorent \emph{pas} ce théorème : nous traitons $F$
 elle-même, et nous n'avons pas établi l'hypothèse de séparateur pour $F^\star$ ni pour les
 feuilles.
\end{itemize}
Le Tableau~\ref{tab:compare} de la Section~\ref{sec:results} résume la comparaison.

\paragraph{Outils.} Le fait que le déterminant d'une matrice polynomiale s'annule en un
point à un ordre au moins égal à la dimension du noyau en ce point est classique ; il
découle de la forme de Smith locale~\cite{GLR82} et constitue un analogue matriciel des
lemmes de multiplicité de la méthode polynomiale~\cite{DKSS13}. Nous incluons la courte
démonstration. L'anticoncentration des sommes de variables aléatoires indépendantes remonte
à Littlewood et Offord~\cite{LO43}, Erd\H{o}s~\cite{Erd45}, Kolmogorov~\cite{Kol58} et
Rogozin~\cite{Rog61}. La minoration $\E|Z|\ge(\E Z^2)^{3/2}/(\E Z^4)^{1/2}$ que nous
utilisons est la méthode du quatrième moment, qui remonte à Littlewood~\cite{Lit30} ; voir
Berger~\cite{Ber97}. Les séparateurs équilibrés ont été introduits par Lipton et
Tarjan~\cite{LT79} ; leur relation avec la largeur arborescente est due à Robertson et
Seymour~\cite{RS86}, voir aussi~\cite{Bod93}. L'expansion des formules CNF aléatoires est la
source standard de bornes inférieures en complexité des preuves~\cite{CS88,BSW01,BSI10} ;
comme dans~\cite{DL}, nous n'utilisons que deux propriétés élémentaires des 3-CNF
aléatoires, démontrées à l'Annexe~\ref{app:random}.

\paragraph{Barrières pour des méthodes.} Comme ceux de~\cite{DL}, nos résultats sont des
barrières pour une méthode dans un modèle restreint, dans l'esprit des limites de la méthode
laser pour la multiplication de matrices~\cite{AFL15,AVW18}. Ce ne sont pas des bornes
inférieures pour un problème.

\paragraph{Nouveauté.} Nous ne connaissons pas de démonstration antérieure des
Théorèmes~\ref{thm:uniform} et~\ref{thm:general}. Cette appréciation repose sur des
recherches effectuées avec des outils d'IA, et non sur un état de l'art systématique, et la
source~\cite{DL} n'a elle-même que quelques semaines. Les corrections sont les bienvenues.

\section{Cadre et notations}\label{sec:setting}

Les Sections~\ref{sec:def}--\ref{sec:question} reproduisent des éléments de~\cite{DL} ; les
passages en retrait sont des citations\footnote{Note de traduction : les citations sont
reproduites en anglais, langue de la source ; chacune est suivie d'une traduction française,
en italique pour les passages longs et en note pour les passages courts.}. La
Section~\ref{sec:ournot} introduit nos propres notations.

\subsection{Schémas d'interpolation par produit de clauses}\label{sec:def}

Chaque clause $e$ possède un polynôme de satisfaction $\varphi_e$, égal à $1$ lorsque la
clause est satisfaite et à $0$ sinon, de sorte que
$\#F=\sum_{x\in\{0,1\}^n}\prod_e\varphi_e(x)$ \cite[(1)]{DL}. La somme déformée est
\begin{equation}\tag{DL.2}\label{eq:H}
H(T)=\sum_{x\in\{0,1\}^n}\ \prod_e\varphi_e\bigl(\langle P_{e,v}(T),x\rangle\bigr)_{v\in e},
\end{equation}
qui est l'équation~(2) de~\cite{DL}. Ce qui suit est cité de \cite[Section~5]{DL}, leur
équation~(3) étant renumérotée~\eqref{eq:three}.

\begin{dlquote}
Let $F$ be a CNF over $x_1,\dots,x_n$ with clauses of at most $k$ literals, $I_e$ the variable
set of clause $e$, and $O=\{(e,v):v\in I_e\}$ its occurrences, $M=|O|$. Let $b$ bound the
number of occurrences of a variable. The incidence graph $G_F$ joins each clause to its
variables; the primal graph $\mathrm{Pr}(F)$ joins two variables that share a clause.

\textbf{Definition 5.1} (Clause-product interpolation scheme). A scheme for $F$ consists of
the following.
\begin{itemize}
\item[(a)] \emph{Deformation.} For each occurrence $o=(e,v)$ a vector $P_o(T)\in\Q[T]^n$ of
 degree at most $R$ with $P_o(0)=e_v$; $H(T)$ as in (2). Its own coefficient is
 $c_o(T)=(P_o(T))_v$, a polynomial of degree at most $R$ with $c_o(0)=1$. Since
 $\deg\varphi_e=|I_e|$, $\deg H\le MR$.
\item[(b)] \emph{Recovery.} A finite set $X\subset\Q\setminus\{0\}$, orders $q_t\ge0$ for
 $t\in X$, and a linear functional $\Lambda$ on Hermite data
 $(f^{(j)}(t))_{t\in X,\,j\le q_t}$ with $\Lambda(f)=f(0)$ for every $f\in\Q[T]$ of degree at
 most $MR$.
\item[(c)] \emph{Designs.} For each $t\in X$: a partition of $[n]$ into groups
 $S_{t,1},S_{t,2},\dots$; a group $\alpha_t(e)$ for each clause; a subspace
 $\cR_{t,h}\subseteq\Q^{S_{t,h}}$ (the span of the measurement rows) for each group, such
 that for every clause $e$, every $v\in I_e$ and every group $h\ne\alpha_t(e)$,
 \begin{equation}\tag{DL.3}\label{eq:three}
 P_{e,v}(t)\big|_{S_{t,h}}\in\cR_{t,h}.
 \end{equation}
 Write $s_t=\max_h|S_{t,h}|$, $\rho_t=\sum_h\dim\cR_{t,h}$, and $\mu_t$ for the number of
 measured groups ($\cR_{t,h}\neq0$). The interaction graph $\cI_t$ has the groups as
 vertices; for each clause $e$ the groups
 $\{\alpha_t(e)\}\cup\{h\ne\alpha_t(e):P_{e,v}(t)|_{S_{t,h}}\ne0\text{ for some }v\in I_e\}$
 form a clique. Let $w_t=\tw(\cI_t)$.
\end{itemize}
The clause-product form of the preprint's interpolation algorithm (its Sections 2.1--2.8) is
an instance, with $\cR_{t,h}$ the row space of its integer measurement matrix $C_{t,h}$ and
$\Lambda$ Hermite interpolation with $q_t=\lceil 3nR/|X|\rceil$ (Section 2.5 of [the
preprint]). We say that $o=(e,v)$ loses its own variable at $t$ if $c_o(t)=0$.
\end{dlquote}

\begin{dltrad}
Soit $F$ une CNF sur $x_1,\dots,x_n$ dont les clauses ont au plus $k$ littéraux, $I_e$
l'ensemble des variables de la clause $e$, et $O=\{(e,v):v\in I_e\}$ ses occurrences,
$M=|O|$. Soit $b$ une borne sur le nombre d'occurrences d'une variable. Le graphe
d'incidence $G_F$ relie chaque clause à ses variables ; le graphe primal $\mathrm{Pr}(F)$
relie deux variables qui partagent une clause.

\textbf{Définition 5.1} (Schéma d'interpolation par produit de clauses). Un schéma pour $F$
est constitué des données suivantes.
\begin{itemize}
\item[(a)] \emph{Déformation.} Pour chaque occurrence $o=(e,v)$, un vecteur
 $P_o(T)\in\Q[T]^n$ de degré au plus $R$ avec $P_o(0)=e_v$ ; $H(T)$ comme en (2). Son
 coefficient propre est $c_o(T)=(P_o(T))_v$, un polynôme de degré au plus $R$ avec
 $c_o(0)=1$. Comme $\deg\varphi_e=|I_e|$, $\deg H\le MR$.
\item[(b)] \emph{Reconstruction.} Un ensemble fini $X\subset\Q\setminus\{0\}$, des ordres
 $q_t\ge0$ pour $t\in X$, et une fonctionnelle linéaire $\Lambda$ sur les données de Hermite
 $(f^{(j)}(t))_{t\in X,\,j\le q_t}$ telle que $\Lambda(f)=f(0)$ pour tout $f\in\Q[T]$ de
 degré au plus $MR$.
\item[(c)] \emph{Designs.} Pour chaque $t\in X$ : une partition de $[n]$ en groupes
 $S_{t,1},S_{t,2},\dots$ ; un groupe $\alpha_t(e)$ pour chaque clause ; un sous-espace
 $\cR_{t,h}\subseteq\Q^{S_{t,h}}$ (l'espace engendré par les lignes de mesure) pour chaque
 groupe, tels que pour toute clause $e$, tout $v\in I_e$ et tout groupe $h\ne\alpha_t(e)$,
 on ait $P_{e,v}(t)\big|_{S_{t,h}}\in\cR_{t,h}$ \textup{(DL.3)}.
 On note $s_t=\max_h|S_{t,h}|$, $\rho_t=\sum_h\dim\cR_{t,h}$, et $\mu_t$ le nombre de
 groupes mesurés ($\cR_{t,h}\neq0$). Le graphe d'interaction $\cI_t$ a pour sommets les
 groupes ; pour chaque clause $e$, les groupes
 $\{\alpha_t(e)\}\cup\{h\ne\alpha_t(e):P_{e,v}(t)|_{S_{t,h}}\ne0\text{ pour un certain }v\in I_e\}$
 forment une clique. Soit $w_t=\tw(\cI_t)$.
\end{itemize}
La forme produit de clauses de l'algorithme d'interpolation de la prépublication (ses
Sections 2.1--2.8) en est une instance, avec pour $\cR_{t,h}$ l'espace des lignes de sa
matrice de mesure entière $C_{t,h}$ et pour $\Lambda$ l'interpolation de Hermite avec
$q_t=\lceil 3nR/|X|\rceil$ (Section 2.5 de [la prépublication]). On dit que $o=(e,v)$ perd
sa propre variable en $t$ si $c_o(t)=0$.
\end{dltrad}

Ici $e_v$ est le $v$-ième vecteur unité de $\Q^n$, et $P_{e,v}(t)|_{S_{t,h}}$ est la
restriction du vecteur $P_{e,v}(t)\in\Q^n$ aux coordonnées de $S_{t,h}$. Nous supposons
toujours $n\ge1$. Nos résultats n'utilisent que les parties suivantes de la définition : les
vecteurs $P_o$ avec $\deg P_o\le R$ et $P_o(0)=e_v$ ; l'exactitude de $\Lambda$ sur tous les
polynômes de degré au plus $MR$ ; et la condition~\eqref{eq:three}. Le polynôme $H$, les
signes des littéraux et la largeur des clauses ne jouent aucun rôle dans les démonstrations.

\subsection{Modèles de coût}\label{sec:cost}

Ce qui suit est cité de \cite[Section~5.2]{DL}.

\begin{dlquote}
We charge each point $t$ at least
\[
\kappa(t)\ \ge\ \max\Bigl(2^{s_t},\ 2^{w_t},\ \tbinom{n}{\min(q_t,\lfloor n/2\rfloor)}\Bigr)
\qquad\text{(DP-type evaluation)},
\]
and, for the preprint's product-type evaluation, additionally $\kappa(t)\ge2^{\rho_t}$. A
scheme runs in time at most $T$ in this model if $\max_{t\in X}\kappa(t)\le T$. All our lower
bounds hold in this model; they are statements about running time only to the extent that the
model's charges are paid.
\end{dlquote}

\begin{dltrad}
Nous facturons à chaque point $t$ au moins
$\kappa(t)\ge\max\bigl(2^{s_t},\,\allowbreak 2^{w_t},\,\allowbreak\tbinom{n}{\min(q_t,\lfloor n/2\rfloor)}\bigr)$
(évaluation de type PD), et, pour l'évaluation de type produit de la prépublication, en
outre $\kappa(t)\ge2^{\rho_t}$. Un schéma s'exécute en temps au plus $T$ dans ce modèle si
$\max_{t\in X}\kappa(t)\le T$. Toutes nos bornes inférieures valent dans ce modèle ; ce ne
sont des énoncés sur le temps d'exécution que dans la mesure où les charges du modèle sont
payées.
\end{dltrad}

Dias et Lima justifient ensuite ces charges. Pour l'évaluation de la prépublication, qui
énumère les affectations de chaque groupe, tous les uplets de lectures de mesure et les
données de Taylor, les charges sont payées. Ils montrent que $w_t\le\mu_t\le\rho_t$, de sorte
que \eng{``the product-type model implies the DP-type one''}\footnote{Traduction : \og le
modèle de type produit implique celui de type PD\fg{}.}. Au sujet du modèle de type PD, ils
écrivent :

\begin{dlquote}
The DP-type model is meant for an evaluation that still counts inside each group by
enumeration, as the preprint does, but sums over measurement readings by dynamic programming
instead of enumerating all tuples. [\dots] Any elimination order over the readings creates a
table indexed by at least $\tw(G_{\mathcal M})+1\ge w_t$ readings, and a nonzero measurement
row takes at least two values on Boolean inputs, so it costs at least $2^{w_t}$. An evaluator
that also counted inside groups by dynamic programming would not pay $2^{s_t}$, and our model
does not cover it.
\end{dlquote}

\begin{dltrad}
Le modèle de type PD est conçu pour une évaluation qui compte encore à l'intérieur de chaque
groupe par énumération, comme le fait la prépublication, mais qui somme sur les lectures de
mesure par programmation dynamique au lieu d'énumérer tous les uplets. [\dots] Tout ordre
d'élimination sur les lectures crée une table indexée par au moins
$\tw(G_{\mathcal M})+1\ge w_t$ lectures, et une ligne de mesure non nulle prend au moins deux
valeurs sur des entrées booléennes, donc cela coûte au moins $2^{w_t}$. Un évaluateur qui
compterait aussi à l'intérieur des groupes par programmation dynamique ne paierait pas
$2^{s_t}$, et notre modèle ne le couvre pas.
\end{dltrad}

Ici $G_{\mathcal M}$ est un graphe sur les groupes mesurés défini dans~\cite{DL} (le torse
de $\cI_t$ sur les groupes mesurés) ; nous n'en avons pas besoin.

On appelle \emph{charge de type PD} d'un schéma le nombre
$\max_{t\in X}\max\bigl(2^{s_t},2^{w_t},\binom{n}{\min(q_t,\lfloor n/2\rfloor)}\bigr)$, et
\emph{charge de type produit} le maximum de la charge de type PD et de
$\max_{t\in X}2^{\rho_t}$. \og Un schéma coûte au moins $C$ dans le modèle de type
produit\fg{} signifie que sa charge de type produit est au moins $C$. Toutes les réserves
de~\cite{DL} s'appliquent à nos énoncés : ils bornent les charges du modèle, et un
évaluateur qui ne paie pas $2^{s_t}$ ou $2^{\rho_t}$ n'est pas couvert.

\subsection{La question de la Section 9 de Dias et Lima}\label{sec:question}

Ce qui suit est cité de \cite[Section~9, \eng{``The gap''}]{DL}\footnote{Traduction :
\og L'écart\fg{}.} ; dans ce passage, $\rho(n)$ désigne l'exposant
$n(\log\log n)^2/\sqrt{\log n}$ de la prépublication, qui est sans rapport avec $\rho_t$.

\begin{dlquote}
On expanding inputs of bounded density, the clause-product scheme is confined between
$2^{\Omega(\sqrt n)}$ (Theorem 8.3) and the preprint's $2^{O(\rho(n))}$, in the cost model.
Our lower-bound technique cannot exceed $2^{\tilde O(\sqrt n)}$: its only constraint is
$(w+1)s+q\gtrsim\tw$, which $s=w=q=\sqrt n$ satisfies. Narrowing the gap needs one of two new
ingredients. From above, a design on an expander with $\sqrt n$ groups, $O(\sqrt n)$
interaction width, Taylor order $O(\sqrt n)$ and $O(1)$ or polylogarithmically many rows per
measured group would bring the scheme to $2^{\tilde O(\sqrt n)}$. From below, one must show
that measurements cannot absorb the cut. A measured group of dimension one can absorb two
foreign occurrences $(e,v),(e',v')$ with $v\ne v'$ at $t$ only if the $2\times2$ minor
$c_{e,v}(t)c_{e',v'}(t)-p_{e,v,v'}(t)p_{e',v',v}(t)$ vanishes; that polynomial equals $1$ at
$T=0$, so each such pair is absorbed at no more than $2R$ points. (Occurrences of the same
variable can be absorbed together at no cost.) Turning this into a bound for groups with
several measurement rows is open.
\end{dlquote}

\begin{dltrad}
Sur les entrées à expansion de densité bornée, le schéma par produit de clauses est confiné
entre $2^{\Omega(\sqrt n)}$ (Théorème 8.3) et le $2^{O(\rho(n))}$ de la prépublication, dans
le modèle de coût. Notre technique de borne inférieure ne peut pas dépasser
$2^{\tilde O(\sqrt n)}$ : sa seule contrainte est $(w+1)s+q\gtrsim\tw$, que $s=w=q=\sqrt n$
satisfait. Réduire l'écart demande l'un de deux ingrédients nouveaux. Par le haut, un design
sur un graphe expanseur avec $\sqrt n$ groupes, une largeur d'interaction $O(\sqrt n)$, un
ordre de Taylor $O(\sqrt n)$ et $O(1)$ lignes, ou un nombre polylogarithmique de lignes, par
groupe mesuré amènerait le schéma à $2^{\tilde O(\sqrt n)}$. Par le bas, il faut montrer que
les mesures ne peuvent pas absorber la coupe. Un groupe mesuré de dimension un ne peut
absorber deux occurrences étrangères $(e,v),(e',v')$ avec $v\ne v'$ en $t$ que si le mineur
$2\times2$ $c_{e,v}(t)c_{e',v'}(t)-p_{e,v,v'}(t)p_{e',v',v}(t)$ s'annule ; ce polynôme vaut
$1$ en $T=0$, donc chaque paire de ce type est absorbée en au plus $2R$ points. (Des
occurrences de la même variable peuvent être absorbées ensemble sans aucun coût.)
Transformer cela en une borne pour des groupes ayant plusieurs lignes de mesure est un
problème ouvert.
\end{dltrad}

Le symbole $p_{e,v,v'}$ n'est pas défini dans~\cite{DL} ; nous le lisons comme la coordonnée
$(P_{e,v}(T))_{v'}$. Trois remarques sur notre lecture de ce passage.

\emph{(1) Il énonce une direction, et non une question formelle.} Nous interprétons la cible
\eng{``from above''}\footnote{Traduction : \og par le haut\fg{}.} comme un schéma dans
lequel $s_t$, $q_t$ et $\rho_t$ sont tous au plus égaux à $\sqrt n\log^{\ell}n$ en chaque
point, pour un certain $\ell$ fixé : avec $\sqrt n$ groupes et un nombre polylogarithmique
de lignes par groupe, $\rho_t=\tilde O(\sqrt n)$, et le coût $2^{\tilde O(\sqrt n)}$ exige
$s_t=\tilde O(\sqrt n)$ dans les deux modèles de coût, parce que $2^{s_t}$ est facturé. Ce
dernier point dépend entièrement de la charge $2^{s_t}$, donc de la réserve citée à la
Section~\ref{sec:cost}.

\emph{(2) Le mineur $2\times2$ est une observation par paires.} Si $v,v'\in S_{t,h}$,
$\dim\cR_{t,h}=1$ et si les deux clauses $e,e'$ sont affectées à d'autres groupes, alors
\eqref{eq:three} place les restrictions de $P_{e,v}(t)$ et de $P_{e',v'}(t)$ à $S_{t,h}$ sur
une même droite, donc la matrice $2\times2$ de leurs coordonnées $v,v'$ est singulière.
Dans~\cite{DL}, cela n'est pas transformé en une borne, même pour une ligne par groupe.

\emph{(3) La phrase entre parenthèses explique notre construction.} Des occurrences de la
même variable peuvent être absorbées ensemble, donc une seule occurrence par variable peut
être utilisée dans un déterminant. C'est pourquoi nous travaillons avec des
\emph{applications témoins}, qui sélectionnent une occurrence de chaque variable. La matrice
$2\times2$ de (2) est le mineur de notre matrice témoin sur les lignes et colonnes $v,v'$,
pour une application témoin qui sélectionne $(e,v)$ et $(e',v')$.

\subsection{Nos notations}\label{sec:ournot}

Le Tableau~\ref{tab:notation} donne la liste des notations. Nous notons $\Gr(F)$ le graphe
primal (noté $\mathrm{Pr}(F)$ dans~\cite{DL}), $\PP$ la probabilité et $\E$ l'espérance, et
$(y)_+=\max(y,0)$. Pour une variable $v$, soit $O_v=\{(e,v):v\in I_e\}$ l'ensemble de ses
occurrences et $d_v=|O_v|$. Ainsi $M=\sum_vd_v$. Soit $\nocc$ le nombre de variables avec
$d_v\ge1$ ; on a clairement $\nocc\le\min(n,M)$.

Fixons un schéma et un point $t\in X$. Lorsque $t$ est fixé, nous écrivons $S_h=S_{t,h}$,
$a_h=|S_h|$ et $r_h=\dim\cR_{t,h}$, de sorte que $\sum_ha_h=n$, $\sum_hr_h=\rho_t$ et
$a_h\le s_t$.

\begin{definition}[Applications témoins, matrice témoin]\label{def:witness}
Une \emph{application témoin} $\omega$ associe à chaque variable $v$ avec $d_v\ge1$ l'une de
ses occurrences $\omega(v)\in O_v$. La \emph{matrice témoin} $\Phi^\omega(T)$ est la matrice
$n\times n$ sur $\Q[T]$ dont la ligne $v$ est $P_{\omega(v)}(T)$ si $d_v\ge1$, et le vecteur
constant $e_v$ si $d_v=0$.
\end{definition}

\begin{definition}[Groupes d'attache, déficit]\label{def:def}
Soit $t\in X$ et soit $\omega$ une application témoin. Le \emph{groupe d'attache}
$h^\omega_t(v)$ d'une variable $v$ est $\alpha_t(e)$ si $\omega(v)=(e,v)$ ; si $d_v=0$,
c'est le groupe qui contient $v$. On pose $I^\omega_{t,h}=\{v:h^\omega_t(v)=h\}$ et
$N^\omega_{t,h}=|I^\omega_{t,h}|$. Le \emph{déficit} en $t$ est
\[
\dfc^\omega_t=\sum_h\bigl(a_h-N^\omega_{t,h}\bigr)_+ .
\]
\end{definition}

Chaque variable a exactement un groupe d'attache, donc $\sum_hN^\omega_{t,h}=n$.

\begin{definition}[Variables scindées]\label{def:split}
Une variable $v$ est \emph{scindée en $t$} si elle apparaît dans deux clauses $e,e'$ avec
$\alpha_t(e)\ne\alpha_t(e')$. Soit $\Sigma_t$ l'ensemble des variables scindées ; tout
$v\in\Sigma_t$ vérifie $d_v\ge2$. Une variable $v\notin\Sigma_t$ a un groupe d'attache
$h_t(v)$ qui ne dépend pas de $\omega$ : si $d_v\ge1$, c'est la valeur commune de
$\alpha_t(e)$ sur les clauses $e\ni v$, et si $d_v=0$, c'est le groupe qui contient $v$. On
pose $J_{t,h}=\{v\notin\Sigma_t:h_t(v)=h\}$. Les ensembles $J_{t,h}$ forment une partition
de $[n]\setminus\Sigma_t$.
\end{definition}

\begin{definition}[Séparateurs]\label{def:sep}
Un ensemble $Y\subseteq[n]$ est un \emph{$2/3$-séparateur} de $F$ si chaque composante
connexe de $\Gr(F)-Y$ a au plus $2n/3$ sommets. Soit $\sep(F)$ la taille minimale d'un
$2/3$-séparateur. Soit $\theta(d)=2(d-1)/d^2$, et
\[
\sepw(F)=\min\Bigl\{\sum_{v\in Y}\theta(d_v)\ :\ Y\text{ un }2/3\text{-séparateur de }F
\text{ avec }d_v\ge2\text{ pour tout }v\in Y\Bigr\},
\]
avec $\sepw(F)=+\infty$ s'il n'existe aucun tel $Y$, et la convention $\min(+\infty,y)=y$.
\end{definition}

La fonction $\theta$ est décroissante sur $d\ge2$, car
$\frac{d}{dd}\frac{d-1}{d^2}=\frac{2-d}{d^3}\le0$. Par conséquent, si $d_v\le b$ pour tout
$v$ avec $b\ge2$,
\begin{equation}\label{eq:sepw}
\sepw(F)\ \ge\ \theta(b)\,\sep(F)=\frac{2(b-1)}{b^2}\,\sep(F).
\end{equation}

\paragraph{Applications témoins aléatoires.} Les applications témoins forment l'ensemble
produit $\prod_{v:d_v\ge1}O_v$. Une \emph{application témoin uniformément aléatoire} est un
élément uniforme de cet ensemble fini ; de manière équivalente, les choix $\omega(v)$ sont
indépendants et $\omega(v)$ est uniforme sur $O_v$. Pour $t\in X$, une variable $v$ et un
groupe $h$, soit
\[
\pi_{t,v,h}=\PP\bigl(h^\omega_t(v)=h\bigr)=
\begin{cases}
|\{e\ni v:\alpha_t(e)=h\}|/d_v&\text{si }d_v\ge1,\\
1\text{ si }v\in S_{t,h},\ 0\text{ sinon}&\text{si }d_v=0,
\end{cases}
\]
et $\Nbar_{t,h}=\E N^\omega_{t,h}=\sum_v\pi_{t,v,h}$. Comme $\sum_h\pi_{t,v,h}=1$, on a
$\sum_h\Nbar_{t,h}=n$. Nous omettons l'indice $t$ lorsque le point est fixé. Les deux
constantes de l'article sont
\[
c_0=\frac{1-e^{-2}}{4}=0.21616\ldots,\qquad
c_b=\frac{c_0}{2\sqrt2}\,\theta(b)=\frac{c_0\,(b-1)}{\sqrt2\,b^2}\quad(b\ge2).
\]

\begin{table}[t]
\centering
\small
\begin{tabular}{@{}>{\raggedright\arraybackslash}p{3.1cm}>{\raggedright\arraybackslash}p{9.6cm}l@{}}
\toprule
Symbole & Signification & Source\\
\midrule
$n$, $M$, $b$ & nombre de variables, d'occurrences ; borne sur le nombre d'occurrences par variable & \cite{DL}\\
$O_v$, $d_v$, $\nocc$ & occurrences de $v$ ; leur nombre ; nombre de $v$ avec $d_v\ge1$ & ici\\
$\Gr(F)$ & graphe primal ($\mathrm{Pr}(F)$ dans~\cite{DL}) & \cite{DL}\\
$P_o(T)$, $R$, $c_o(T)$ & vecteur de déformation, borne sur le degré, coefficient propre & \cite{DL}\\
$X$, $q_t$, $\Lambda$ & points d'évaluation, ordres de Taylor, fonctionnelle de reconstruction & \cite{DL}\\
$S_{t,h}$, $\alpha_t$, $\cR_{t,h}$ & groupes, affectation des clauses, sous-espaces de mesure & \cite{DL}\\
$a_h$, $r_h$, $g_t$ & $|S_{t,h}|$, $\dim\cR_{t,h}$, nombre de groupes en $t$ & ici\\
$s_t$, $\rho_t$, $\mu_t$, $w_t$ & $\max_ha_h$, $\sum_hr_h$, nombre de groupes mesurés, $\tw(\cI_t)$ & \cite{DL}\\
$\kappa(t)$ & coût facturé au point $t$ & \cite{DL}\\
$\omega$, $\Phi^\omega(T)$ & application témoin, matrice témoin & Déf.~\ref{def:witness}\\
$h^\omega_t(v)$, $I^\omega_{t,h}$, $N^\omega_{t,h}$ & groupe d'attache ; variables de groupe d'attache $h$ ; leur nombre & Déf.~\ref{def:def}\\
$\dfc^\omega_t$ & déficit $\sum_h(a_h-N^\omega_{t,h})_+$ & Déf.~\ref{def:def}\\
$\pi_{t,v,h}$, $\Nbar_{t,h}$ & $\PP(h^\omega_t(v)=h)$, $\E N^\omega_{t,h}$ pour $\omega$ uniforme & ici\\
$\Sigma_t$, $h_t(v)$, $J_{t,h}$ & variables scindées ; groupe d'attache d'une variable non scindée ; ses fibres & Déf.~\ref{def:split}\\
$\sep(F)$, $\theta(d)$, $\sepw(F)$ & nombre de séparation ; poids $2(d-1)/d^2$ ; nombre de séparation pondéré & Déf.~\ref{def:sep}\\
$c_0$, $c_b$ & $(1-e^{-2})/4$ ; $c_0(b-1)/(\sqrt2b^2)$ & ici\\
\bottomrule
\end{tabular}
\caption{Notations. Le nombre de groupes mesurés $\mu_t$ de~\cite{DL} n'apparaît qu'à la
Section~\ref{sec:related} et à la Section~\ref{sec:dp}.}\label{tab:notation}
\end{table}

\section{Résultats principaux}\label{sec:results}

Dans toute la suite, $F$ est une formule CNF sur $n\ge1$ variables et un \emph{schéma} est un
schéma d'interpolation par produit de clauses pour $F$ au sens de la Définition~5.1
de~\cite{DL} (Section~\ref{sec:def}). Aucun modèle de coût n'est utilisé avant le
Corollaire~\ref{cor:main}.

\begin{lemma}[Comptage de Hermite]\label{lem:hermite}
Dans tout schéma, $\sum_{t\in X}(q_t+1)\ge MR+1$. En particulier $X\ne\emptyset$.
\end{lemma}

\begin{lemma}[Déterminant témoin]\label{lem:det}
Pour tout schéma et toute application témoin $\omega$,
\[
\sum_{t\in X}\bigl(\dfc^\omega_t-\rho_t\bigr)_+\ \le\ \sum_{t\in X}\dim\ker\Phi^\omega(t)
\ \le\ \deg\det\Phi^\omega\ \le\ \nocc R\ \le\ \min(n,M)\,R .
\]
Plus précisément, pour tout $t\in X$,
$\dim\ker\Phi^\omega(t)\ge\sum_h\bigl(a_h-r_h-N^\omega_{t,h}\bigr)_+\ge\dfc^\omega_t-\rho_t$.
\end{lemma}

La Figure~\ref{fig:blocks} montre la structure qui sous-tend le Lemme~\ref{lem:det}.

\begin{figure}[t]
\[
\Phi^\omega(t)\;=\;
\begin{array}{r|ccc}
 & \text{colonnes } S_{1} & \text{colonnes } S_{2} & \text{colonnes } S_{3}\\ \hline
\text{lignes } I^\omega_{t,1} & \text{libre} & \in\cR_{t,2} & \in\cR_{t,3}\\
\text{lignes } I^\omega_{t,2} & \in\cR_{t,1} & \text{libre} & \in\cR_{t,3}\\
\text{lignes } I^\omega_{t,3} & \in\cR_{t,1} & \in\cR_{t,2} & \text{libre}
\end{array}
\]
\caption{Structure par blocs de la matrice témoin en un point $t$, pour trois groupes. Les
lignes sont indexées par les variables, triées par groupe d'attache ; les colonnes sont
triées par groupe. D'après la condition~\eqref{eq:three}, une ligne dont le groupe d'attache
n'est pas $h$ a sa restriction aux colonnes $S_h$ dans $\cR_{t,h}$. Le bloc de colonnes
$S_h$ est donc de rang au plus $\min(a_h,\,N^\omega_{t,h}+r_h)$. Noter que $N^\omega_{t,h}$,
le nombre de lignes de groupe d'attache $h$, n'est pas nécessairement égal à
$a_h=|S_h|$.}\label{fig:blocks}
\end{figure}

\begin{lemma}[Anticoncentration]\label{lem:anti}
Soit $B$ une somme d'un nombre fini de variables aléatoires de Bernoulli indépendantes, de
variance $V$, et soit $a\in\mathbb R$. Alors $\E|B-a|\ge c_0\min(V,\sqrt V)$, où
$c_0=(1-e^{-2})/4$.
\end{lemma}

\begin{lemma}[Espérance du déficit]\label{lem:expdef}
Soit $\omega$ une application témoin uniformément aléatoire. Pour tout schéma et tout
$t\in X$,
\[
\E\,\dfc^\omega_t\ \ge\ \max\Bigl(\frac12\sum_h\bigl|\Nbar_{t,h}-a_h\bigr|,\ \
\frac{c_0}{2\sqrt{2s_t}}\sum_{v\in\Sigma_t}\theta(d_v)\Bigr).
\]
\end{lemma}

\begin{lemma}[Les variables scindées forment un séparateur]\label{lem:split}
Pour tout schéma et tout $t\in X$, l'une au moins des assertions suivantes est vraie :
\begin{enumerate}
\item[\textup{(1)}] $|J_{t,h}|\le 2n/3$ pour tout groupe $h$. Alors $\Sigma_t$ est un
 $2/3$-séparateur de $F$ ; en particulier $|\Sigma_t|\ge\sep(F)$ et
 $\sum_{v\in\Sigma_t}\theta(d_v)\ge\sepw(F)$.
\item[\textup{(2)}] $|J_{t,h_0}|>2n/3$ pour un certain groupe $h_0$. Alors
 $\frac12\sum_h|\Nbar_{t,h}-a_h|>\frac{2n}{3}-s_t$.
\end{enumerate}
\end{lemma}

\begin{theorem}[Déformations uniformes par variable : aucune absorption]\label{thm:uniform}
Supposons que la déformation est uniforme par variable, c'est-à-dire que
$P_{e,v}=P_{e',v}$ pour toutes clauses $e,e'$ contenant $v$. Alors un certain $t\in X$
vérifie
\[
\rho_t+q_t+1\ >\ \min\Bigl(\sep(F),\ \frac{2n}{3}-s_t\Bigr).
\]
\end{theorem}

\begin{theorem}[Déformations générales : au plus un facteur $\sqrt s$]\label{thm:general}
Tout schéma pour $F$ possède un point $t\in X$ tel que
\[
\rho_t+q_t+1\ >\ \min\Bigl(\frac{c_0}{2\sqrt2}\cdot\frac{\sepw(F)}{\sqrt{s_t}},\ \
\frac{2n}{3}-s_t\Bigr).
\]
Ceci vaut pour des déformations dépendant de la clause, des partitions, des affectations
$\alpha_t$ et des sous-espaces $\cR_{t,h}$ arbitraires. Si chaque variable apparaît au plus
$b$ fois, $b\ge2$, alors un certain $t\in X$ vérifie
\[
\rho_t+q_t+1\ >\ \min\Bigl(c_b\,\frac{\sep(F)}{\sqrt{s_t}},\ \frac{2n}{3}-s_t\Bigr),
\qquad c_b=\frac{c_0\,(b-1)}{\sqrt2\,b^2},\quad c_3=\frac{1-e^{-2}}{18\sqrt2}=0.03396\ldots
\]
\end{theorem}

La constante $c_0$ n'est pas optimisée (Remarque~\ref{rem:const}).

\begin{lemma}[3-CNF aléatoires]\label{lem:random}
Fixons $\Delta>3$. Soit $F$ une 3-CNF aléatoire sur $n$ variables avec $m=\Delta n$ clauses
indépendantes (on suppose que $\Delta n$ est un entier), chacune portant sur trois variables
distinctes uniformément aléatoires avec des signes aléatoires. Il existe des constantes
$0<\eta'_\Delta\le\eta_\Delta\le1/9$ telles que, avec une probabilité tendant vers $1$
lorsque $n\to\infty$, $\sep(F)\ge\eta_\Delta n$ et $\sepw(F)\ge\eta'_\Delta n$.
\end{lemma}

C'est le modèle de \cite[Annexe~A]{DL} ; la démonstration adapte celle de
\cite[Lemme~8.2]{DL}.

\begin{corollary}\label{cor:main}
Soit $F$ appartenant à l'une des familles suivantes.
\begin{itemize}
\item[\textup{(i)}] Chaque variable apparaît au plus $b$ fois et $\sep(F)\ge cn$, pour
 $b\ge2$ et $c>0$ fixés. On pose $\gamma=c_bc$.
\item[\textup{(ii)}] $F$ est une 3-CNF aléatoire comme au Lemme~\ref{lem:random}, sur
 l'événement de ce lemme (qui est de probabilité $1-o(1)$). On pose
 $\gamma=\frac{c_0}{2\sqrt2}\eta'_\Delta$ et $c=\eta_\Delta$.
\end{itemize}
Dans les deux cas, $\gamma\le1/6$ et $c\le1$. Il existe $n_0$, ne dépendant que de $\gamma$
et, dans la partie~1, de $\ell$, tel que pour $n\ge n_0$ les assertions suivantes valent pour
tout schéma pour $F$.
\begin{enumerate}
\item \textup{(Le design de \cite[Section~9]{DL} n'existe pas.)} Soit $\ell\ge0$ fixé.
 Si $s_t\le\sqrt n\log^{\ell}n$ pour tout $t\in X$, alors un certain $t\in X$ vérifie
 \[\rho_t+q_t+1>\gamma\,n^{3/4}/\log^{\ell/2}n .\]
 En particulier, aucun schéma ne vérifie $s_t,\ q_t,\ \rho_t\le\sqrt n\log^{\ell}n$ pour tout
 $t\in X$. Cet énoncé ne fait intervenir aucun modèle de coût.
\item \textup{(Coût de type produit.)} La charge de type produit du schéma est au moins
 $2^{(\gamma n^{2/3}-1)/2}=2^{\Omega(n^{2/3})}$.
\item \textup{(Schémas uniformes par variable.)} Si la déformation est uniforme par
 variable, la charge de type produit est au moins
 $2^{\min(n/6,\ (\min(c,1/2)\,n-1)/2)}=2^{\Omega(n)}$. Dans la
 famille~\textup{(i)}, cette partie n'a pas besoin de la borne $b$ sur les occurrences.
\end{enumerate}
\end{corollary}

Dans les trois parties, le schéma est appliqué à $F$ elle-même. Pour la famille~(ii), cela
signifie : à la formule aléatoire, et non à la formule à occurrences bornées $F^\star$ que
la prépublication en dérive, ni aux feuilles de la sparsification.

\begin{remark}[Designs équilibrés]\label{rem:balanced}
Soit $F$ comme au Corollaire~\ref{cor:main} et soit $t$ le point donné par le
Théorème~\ref{thm:general}. Supposons que, en ce $t$ particulier, $s_t\le n/2$, et que le
design a $g_t$ groupes avec $s_t\le Kn/g_t$ pour un certain $K\ge1$. Alors
$q_t+1\ge\frac{\gamma}{2\sqrt K}\sqrt{ng_t}$, ou bien le nombre moyen de lignes par groupe
vérifie $\rho_t/g_t>\frac{\gamma}{2K}\sqrt{s_t}$. La démonstration se trouve à la
Section~\ref{sec:proofcor}. L'énoncé concerne un seul point $t$, et non tous les points, et
il a besoin de $s_t\le n/2$ pour que le premier terme du minimum du
Théorème~\ref{thm:general} soit le plus petit.
\end{remark}

\begin{table}[t]
\centering
\small
\begin{tabular}{@{}>{\raggedright\arraybackslash}p{2.5cm}>{\raggedright\arraybackslash}p{3.7cm}>{\raggedright\arraybackslash}p{2.1cm}>{\raggedright\arraybackslash}p{6.2cm}@{}}
\toprule
Résultat & Schéma appliqué à & Modèle & Borne\\
\midrule
\cite[Th.~7.1(ii)]{DL} & $F$ quelconque avec $s_t\le 2n/3$ & aucun & $q_t+b\,\mu_ts_t\ge\sep(F)$ en un certain $t$ \newline(groupes mesurés facturés)\\[2pt]
\cite[Th.~8.1]{DL} & $F$ quelconque & type PD & coût $\ge2^{E}$, \newline $E=\min(\sqrt{(\tw+1)/2}-1,\,(\tw+1)/2,\,\lfloor n/2\rfloor)$\\[2pt]
\cite[Th.~8.3]{DL} & $F^\star$ et feuilles de la sparsification, 3-CNF aléatoire, $\Delta>3$ & type PD & coût $\ge2^{\Omega_\Delta(\sqrt n)}$ avec forte probabilité\\[2pt]
Prop.~\ref{prop:dp} & $F$ quelconque & type PD & un schéma trivial a une charge de $2^{\lceil\sqrt n\rceil}$ (borne supérieure)\\[2pt]
Th.~\ref{thm:uniform} & $F$ quelconque, uniforme par variable & aucun & $\rho_t+q_t+1>\min(\sep(F),\,2n/3-s_t)$ \newline en un certain $t$\\[2pt]
Th.~\ref{thm:general} & $F$ quelconque avec $d_v\le b$ & aucun & $\rho_t+q_t+1>\min(c_b\sep(F)/\sqrt{s_t},\,2n/3-s_t)$ \newline en un certain $t$ (lignes facturées)\\[2pt]
Cor.~\ref{cor:main}(2) & $F$ elle-même, familles (i), (ii) & type produit & coût $\ge2^{\Omega(n^{2/3})}$\\[2pt]
Cor.~\ref{cor:main}(3) & $F$ elle-même, uniforme par variable & type produit & coût $\ge2^{\Omega(n)}$\\
\bottomrule
\end{tabular}
\caption{Comparaison avec les bornes de~\cite{DL}. La première ligne est notre lecture de
\cite[Th.~7.1(ii)]{DL} en termes de séparateurs (Section~\ref{sec:related}). Les bornes
\og de type PD\fg{} valent aussi dans le modèle de type produit, qui facture davantage. Les
lignes relatives à~\cite{DL} et celles relatives à cet article concernent des formules
différentes dans le cas du Théorème~8.3 ($F^\star$ et les feuilles là-bas, $F$
ici).}\label{tab:compare}
\end{table}

\section{Démonstrations}\label{sec:proofs}

\subsection{Démonstration du Lemme~\ref{lem:hermite} (comptage de Hermite)}

L'argument est celui de \cite[Lemme~6.2]{DL}, appliqué à un autre polynôme. Supposons
$\sum_{t\in X}(q_t+1)\le MR$, et soit $f(T)=\prod_{t\in X}(T-t)^{q_t+1}$ (donc $f=1$ si
$X=\emptyset$). Alors $\deg f\le MR$. Pour tout $t\in X$ et tout $j\le q_t$, on a
$f^{(j)}(t)=0$, car $(T-t)^{q_t+1}$ divise $f$. Les données de Hermite de $f$ sont donc
toutes nulles et, $\Lambda$ étant linéaire, $\Lambda(f)=0$. Mais
$f(0)=\prod_{t\in X}(-t)^{q_t+1}\ne0$ car $0\notin X$. Ceci contredit l'exactitude
$\Lambda(f)=f(0)$ exigée dans la partie~(b) de la définition. Par conséquent
$\sum_{t\in X}(q_t+1)\ge MR+1\ge1$, et $X\ne\emptyset$. \qed

\subsection{Démonstration du Lemme~\ref{lem:det} (déterminant témoin)}

Fixons un schéma et une application témoin $\omega$, et écrivons $\Phi=\Phi^\omega$.

\paragraph{Étape 1 : $\det\Phi$ est un polynôme non nul de degré au plus $\nocc R$.} Si
$d_v\ge1$, la ligne $v$ de $\Phi(0)$ est $P_{\omega(v)}(0)=e_v$, car $\omega(v)$ est une
occurrence de $v$. Si $d_v=0$, la ligne $v$ est $e_v$ par définition. Donc $\Phi(0)$ est la
matrice identité et $\det\Phi(0)=1$ ; en particulier $\det\Phi\ne0$. D'après la formule de
Leibniz, $\det\Phi$ est une somme signée de produits de $n$ coefficients, un dans chaque
ligne. Les coefficients d'une ligne $v$ avec $d_v\ge1$ sont de degré au plus $R$, et les
coefficients d'une ligne avec $d_v=0$ sont des constantes. Donc chaque produit est de degré
au plus $\nocc R$, et $\deg\det\Phi\le\nocc R\le\min(n,M)R$.

\paragraph{Étape 2 : borne sur le rang en un point.} Fixons $t\in X$ et un groupe $h$. Soit
$B_h$ la matrice $n\times a_h$ formée des colonnes de $\Phi(t)$ indexées par $S_h$. Soit $v$
une variable avec $h^\omega_t(v)\ne h$. Nous affirmons que la ligne $v$ de $B_h$ appartient
à $\cR_{t,h}$.
\begin{itemize}
\item Si $d_v\ge1$ et $\omega(v)=(e,v)$, alors $\alpha_t(e)=h^\omega_t(v)\ne h$, et la ligne
 $v$ de $B_h$ est $P_{e,v}(t)|_{S_h}$, qui appartient à $\cR_{t,h}$
 d'après~\eqref{eq:three}.
\item Si $d_v=0$, alors $h^\omega_t(v)$ est le groupe contenant $v$, donc $v\notin S_h$, et
 la ligne $v$ de $B_h$ est $e_v|_{S_h}=0\in\cR_{t,h}$.
\end{itemize}
Donc l'espace des lignes de $B_h$ est contenu dans la somme de $\cR_{t,h}$ et de l'espace
engendré par les $N^\omega_{t,h}$ lignes indexées par $I^\omega_{t,h}$. Comme $B_h$ a $a_h$
colonnes,
\[
\rank B_h\ \le\ \min\bigl(a_h,\ r_h+N^\omega_{t,h}\bigr).
\]
Les groupes forment une partition des colonnes de $\Phi(t)$, donc l'espace des colonnes de
$\Phi(t)$ est la somme des espaces des colonnes des $B_h$, et
$\rank\Phi(t)\le\sum_h\rank B_h$. Avec $\sum_ha_h=n$, cela donne
\[
\dim\ker\Phi(t)=n-\rank\Phi(t)\ \ge\ \sum_h\Bigl(a_h-\min\bigl(a_h,r_h+N^\omega_{t,h}\bigr)\Bigr)
=\sum_h\bigl(a_h-r_h-N^\omega_{t,h}\bigr)_+ .
\]
Pour $y$ réel et $r\ge0$, on a $(y-r)_+\ge y_+-r$. En appliquant ceci avec
$y=a_h-N^\omega_{t,h}$ et $r=r_h$ et en sommant sur $h$,
\[
\dim\ker\Phi(t)\ \ge\ \sum_h\bigl(a_h-N^\omega_{t,h}\bigr)_+-\sum_hr_h=\dfc^\omega_t-\rho_t .
\]
Comme $\dim\ker\Phi(t)\ge0$, on a aussi $\dim\ker\Phi(t)\ge(\dfc^\omega_t-\rho_t)_+$.

\paragraph{Étape 3 : ordre d'annulation.} Fixons $t\in X$ et soit $k=\dim\ker\Phi(t)$.
Choisissons une matrice inversible $U\in\Q^{n\times n}$ dont les $k$ premières colonnes
forment une base de $\ker\Phi(t)$. La matrice $\Phi(T)U$ est à coefficients dans $\Q[T]$, et
ses $k$ premières colonnes s'annulent en $T=t$. Donc chaque coefficient de ces colonnes est
divisible par $T-t$. En mettant $T-t$ en facteur dans chacune de ces $k$ colonnes, on
obtient $\det(\Phi(T)U)=(T-t)^k\det M(T)$ pour une matrice $M(T)$ sur $\Q[T]$. Comme
$\det(\Phi U)=\det U\cdot\det\Phi$ et que $\det U$ est une constante non nulle,
$(T-t)^{k}$ divise $\det\Phi$.

\paragraph{Étape 4 : conclusion.} Les polynômes $(T-t)^{\dim\ker\Phi(t)}$, $t\in X$, sont
deux à deux premiers entre eux et chacun divise $\det\Phi$. Donc leur produit divise
$\det\Phi$, qui est non nul d'après l'étape~1. En comparant les degrés,
$\sum_{t\in X}\dim\ker\Phi(t)\le\deg\det\Phi\le\nocc R$. Avec l'étape~2, cela démontre le
lemme. \qed

\begin{remark}\label{rem:mn}
Seule l'inégalité $\deg\det\Phi^\omega\le MR$ est utilisée dans la suite. Elle vaut pour
toute formule parce que $\nocc\le M$ ; aucune hypothèse telle que \og toute variable
apparaît\fg{} ou $M\ge n$ n'est nécessaire.
\end{remark}

\subsection{Démonstration du Lemme~\ref{lem:anti} (anticoncentration)}

Écrivons $B=\sum_{v}\beta_v$ avec des $\beta_v\sim\mathrm{Bernoulli}(p_v)$ indépendantes,
de sorte que $V=\sum_vp_v(1-p_v)$. Si $V=0$, il n'y a rien à démontrer. Soit
$B'=\sum_v\beta'_v$ une copie indépendante de $B$. Par l'inégalité triangulaire,
$|B-B'|\le|B-a|+|B'-a|$, donc
\begin{equation}\label{eq:sym}
\E|B-a|\ \ge\ \tfrac12\,\E|D|,\qquad D=B-B'=\sum_v\xi_v,\quad \xi_v=\beta_v-\beta'_v .
\end{equation}
Les $\xi_v$ sont indépendantes, prennent leurs valeurs dans $\{-1,0,1\}$, et sont
symétriques : $\PP(\xi_v=1)=\PP(\xi_v=-1)=p_v(1-p_v)$. Posons
$\tau_v=\PP(\xi_v\ne0)=2p_v(1-p_v)$. Alors $\E\xi_v=0$, $\E\xi_v^2=\E\xi_v^4=\tau_v$, et
$\sum_v\tau_v=2V$. En particulier $\E D^2=\sum_v\tau_v=2V$.

\paragraph{(i) Une borne pour tout $V$.} Soit $K=\{v:\xi_v\ne0\}$. Conditionnellement à
$K$, les valeurs $(\xi_v)_{v\in K}$ sont des signes uniformes indépendants, car les $\xi_v$
sont indépendantes et symétriques. Si $|K|=j\ge1$, alors $\PP(D=0\mid K)$ vaut $0$ pour $j$
impair et $\binom{j}{j/2}2^{-j}\le\frac12$ pour $j$ pair. Donc
$\PP(D\ne0\mid K)\ge\frac12$ dès que $K\ne\emptyset$. De plus
\[
\PP(K\ne\emptyset)=1-\prod_v(1-\tau_v)\ \ge\ 1-e^{-\sum_v\tau_v}=1-e^{-2V}.
\]
Comme $D$ est un entier, $|D|\ge1$ lorsque $D\ne0$, donc
$\E|D|\ge\PP(D\ne0)\ge\frac12(1-e^{-2V})$, et d'après~\eqref{eq:sym},
\begin{equation}\label{eq:smallV}
\E|B-a|\ \ge\ \tfrac14\bigl(1-e^{-2V}\bigr).
\end{equation}

\paragraph{(ii) Une borne pour $V\ge1$.} En développant $D^4$ et en utilisant
l'indépendance et $\E\xi_v=\E\xi_v^3=0$, seuls les termes $\xi_v^4$ et $\xi_u^2\xi_v^2$
($u\ne v$) ont une espérance non nulle, et chaque paire non ordonnée $\{u,v\}$ apparaît
$\binom42=6$ fois :
\[
\E D^4=\sum_v\tau_v+3\sum_{u\ne v}\tau_u\tau_v\ \le\ 2V+3(2V)^2=2V+12V^2 .
\]
D'après l'inégalité de H\"older pour les exposants $3/2$ et $3$, on a
$\E D^2=\E\bigl(|D|^{2/3}|D|^{4/3}\bigr)\le(\E|D|)^{2/3}(\E D^4)^{1/3}$, d'où
\[
\E|D|\ \ge\ \frac{(\E D^2)^{3/2}}{(\E D^4)^{1/2}}\ \ge\ \frac{(2V)^{3/2}}{(2V+12V^2)^{1/2}} .
\]
Pour $V\ge1$, on a $2V+12V^2\le14V^2$, donc $\E|D|\ge\frac{2\sqrt2}{\sqrt{14}}\sqrt V
=\sqrt{4/7}\,\sqrt V$, et d'après~\eqref{eq:sym}
$\E|B-a|\ge\frac{1}{\sqrt7}\sqrt V\ge0.377\sqrt V\ge c_0\sqrt V$.

\paragraph{Conclusion.} La fonction $V\mapsto1-e^{-2V}$ est concave et s'annule en $0$, donc
$(1-e^{-2V})/V$ est décroissante sur $V>0$. Pour $0<V\le1$, cela donne
$1-e^{-2V}\ge(1-e^{-2})V$, et~\eqref{eq:smallV} fournit
$\E|B-a|\ge\frac{1-e^{-2}}{4}V=c_0V$. Pour $V\ge1$, (ii) donne $\E|B-a|\ge c_0\sqrt V$. Comme
$\min(V,\sqrt V)$ vaut $V$ pour $V\le1$ et $\sqrt V$ pour $V\ge1$, le lemme en découle. \qed

\begin{remark}[Sur la constante]\label{rem:const}
La valeur de $c_0$ est imposée par l'étape~(i) en $V=1$ ; l'étape~(ii) donne $0.377$. Nous
n'avons pas cherché à l'optimiser. Dans l'expérience numérique de la
Section~\ref{sec:checks}, le plus petit rapport $\E|B-a|/\min(V,\sqrt V)$ observé était
$0.7397$.
\end{remark}

\subsection{Démonstration du Lemme~\ref{lem:expdef} (espérance du déficit)}

Fixons $t\in X$ et omettons-le dans les notations ; écrivons $s=s_t$ et
$N_h=N^\omega_{t,h}$. Rappelons que $\sum_hN_h=n=\sum_ha_h$. Donc $\sum_h(a_h-N_h)=0$, les
parties positives et négatives des nombres $a_h-N_h$ ont la même somme, et
\[
\dfc^\omega_t=\sum_h(a_h-N_h)_+=\frac12\sum_h|N_h-a_h|,\qquad\text{d'où}\qquad
\E\,\dfc^\omega_t=\frac12\sum_h\E|N_h-a_h| .
\]

\paragraph{Première borne.} Par l'inégalité de Jensen,
$\E|N_h-a_h|\ge|\E N_h-a_h|=|\Nbar_h-a_h|$. En sommant sur $h$, on obtient le premier terme
du maximum.

\paragraph{Structure de $N_h$.} Pour un groupe $h$ fixé,
$N_h=\sum_v\mathbf 1[h^\omega_t(v)=h]$. Le groupe d'attache de $v$ ne dépend que de
$\omega(v)$ (et il est déterministe si $d_v=0$), et les $\omega(v)$ sont indépendants. Donc
$N_h$ est une somme de variables de Bernoulli indépendantes de paramètres $\pi_{v,h}$, de
moyenne $\Nbar_h$ et de variance
\[
V_h=\sum_v\pi_{v,h}(1-\pi_{v,h})\ \le\ \Nbar_h .
\]

\paragraph{Seconde borne, pour un groupe.} Nous affirmons que
\begin{equation}\label{eq:onegroup}
\E|N_h-a_h|\ \ge\ \frac{c_0}{\sqrt{2s}}\,V_h\qquad\text{pour tout groupe }h.
\end{equation}
Noter que $s\ge1$ et $a_h\le s$.

\emph{Cas $\Nbar_h\le2s$.} Alors $V_h\le2s$. Si $V_h\le1$, alors
$\min(V_h,\sqrt{V_h})=V_h\ge V_h/\sqrt{2s}$ car $\sqrt{2s}\ge1$. Si $V_h\ge1$, alors
$\min(V_h,\sqrt{V_h})=\sqrt{V_h}=V_h/\sqrt{V_h}\ge V_h/\sqrt{2s}$. Dans les deux sous-cas,
le Lemme~\ref{lem:anti}, appliqué à $B=N_h$ et $a=a_h$, donne~\eqref{eq:onegroup}.

\emph{Cas $\Nbar_h>2s$.} Alors $a_h\le s<\Nbar_h/2$, donc
$\E|N_h-a_h|\ge\Nbar_h-a_h>\Nbar_h/2\ge V_h/2$. Comme $s\ge1$,
$c_0/\sqrt{2s}\le c_0/\sqrt2<1/2$, et~\eqref{eq:onegroup} en découle.

\paragraph{Somme sur les groupes.} Comme $\sum_h\pi_{v,h}=1$,
\[
\sum_hV_h=\sum_v\Bigl(1-\sum_h\pi_{v,h}^2\Bigr),
\]
et chaque terme de la somme extérieure est positif ou nul. Soit $v\in\Sigma_t$ et
$d=d_v\ge2$. Alors $\pi_{v,h}=k_h/d$, où $k_h$ est le nombre de clauses contenant $v$ qui
sont affectées à $h$ ; les $k_h$ sont des entiers positifs ou nuls de somme $d$, et au moins
deux d'entre eux sont strictement positifs. Soit $k_{h_1}$ l'un des plus grands ; alors
$1\le k_{h_1}\le d-1$ et
$\sum_{h\ne h_1}k_h^2\le\bigl(\sum_{h\ne h_1}k_h\bigr)^2=(d-k_{h_1})^2$. La fonction
$k\mapsto k^2+(d-k)^2$ est convexe, donc sur $1\le k\le d-1$ elle est au plus égale à sa
valeur $(d-1)^2+1$ aux extrémités. Par conséquent $\sum_hk_h^2\le(d-1)^2+1$ et
\[
1-\sum_h\pi_{v,h}^2\ \ge\ 1-\frac{(d-1)^2+1}{d^2}=\frac{2(d-1)}{d^2}=\theta(d_v).
\]
En ne gardant que les variables scindées, $\sum_hV_h\ge\sum_{v\in\Sigma_t}\theta(d_v)$.
Avec~\eqref{eq:onegroup},
\[
\E\,\dfc^\omega_t=\frac12\sum_h\E|N_h-a_h|\ \ge\ \frac{c_0}{2\sqrt{2s}}\sum_hV_h
\ \ge\ \frac{c_0}{2\sqrt{2s}}\sum_{v\in\Sigma_t}\theta(d_v),
\]
ce qui est le second terme du maximum. \qed

\subsection{Démonstration du Lemme~\ref{lem:split} (les variables scindées forment un séparateur)}

Fixons $t\in X$.

\paragraph{Composantes de $\Gr(F)-\Sigma_t$.} Soient $u,v\notin\Sigma_t$ adjacentes dans
$\Gr(F)$. Elles partagent une clause $e$, donc $d_u,d_v\ge1$, et comme aucune des deux n'est
scindée, $h_t(u)=\alpha_t(e)=h_t(v)$ d'après la Définition~\ref{def:split}. Donc $h_t$ est
constante sur chaque composante connexe de $\Gr(F)-\Sigma_t$ : chaque composante est
contenue dans un seul ensemble $J_{t,h}$. (Les variables avec $d_v=0$ sont des sommets
isolés de $\Gr(F)$ ; leur groupe d'attache est le groupe qui les contient, d'après la
convention de la Définition~\ref{def:split}.)

\paragraph{Cas (1).} Si $|J_{t,h}|\le2n/3$ pour tout $h$, chaque composante de
$\Gr(F)-\Sigma_t$ a au plus $2n/3$ sommets, donc $\Sigma_t$ est un $2/3$-séparateur et
$|\Sigma_t|\ge\sep(F)$. Toute variable scindée vérifie $d_v\ge2$, donc $\Sigma_t$ est
admissible dans la définition de $\sepw(F)$ et $\sum_{v\in\Sigma_t}\theta(d_v)\ge\sepw(F)$.

\paragraph{Cas (2).} Supposons $|J_{t,h_0}|>2n/3$. Soit $v\in J_{t,h_0}$. Si $d_v\ge1$,
toutes les clauses contenant $v$ sont affectées à $h_0$, donc $\pi_{v,h_0}=1$. Si $d_v=0$,
alors $v\in S_{t,h_0}$ et $\pi_{v,h_0}=1$ par définition. Donc
$\Nbar_{h_0}\ge|J_{t,h_0}|>2n/3$, tandis que $a_{h_0}\le s_t$. Comme
$\sum_h\Nbar_h=n=\sum_ha_h$, les nombres $\Nbar_h-a_h$ sont de somme nulle, et
\[
\frac12\sum_h\bigl|\Nbar_h-a_h\bigr|=\sum_h\bigl(\Nbar_h-a_h\bigr)_+\ \ge\ \Nbar_{h_0}-a_{h_0}
\ >\ \frac{2n}{3}-s_t .
\]
Les deux cas sont exhaustifs. \qed

\subsection{Démonstration du Théorème~\ref{thm:uniform} (déformations uniformes par variable)}

Supposons $P_{e,v}=P_v$ pour toute clause $e\ni v$. Alors $\Phi^\omega$ ne dépend pas de
l'application témoin : sa ligne $v$ est $P_v(T)$ si $d_v\ge1$ et $e_v$ si $d_v=0$. Notons
$\Phi$ cette matrice. (La déformation substitue alors la forme linéaire
$\langle\Phi_v(T),x\rangle$ à $x_v$ dans chaque clause.)

\paragraph{Borne sur le rang.} Fixons $t\in X$ et un groupe $h$, et soit $B_h$ le bloc des
colonnes $S_h$ de $\Phi(t)$. Nous affirmons que la ligne $v$ de $B_h$ appartient à
$\cR_{t,h}$ pour tout $v\notin J_{t,h}$. Il y a trois cas.
\begin{itemize}
\item[(a)] $v\in\Sigma_t$. Alors $v$ apparaît dans des clauses affectées à deux groupes
 différents, donc elle apparaît dans une clause $e$ avec $\alpha_t(e)\ne h$.
\item[(b)] $v\notin\Sigma_t$, $d_v\ge1$ et $h_t(v)\ne h$. Alors toute clause $e\ni v$
 vérifie $\alpha_t(e)=h_t(v)\ne h$.
\item[(c)] $d_v=0$ et $h_t(v)\ne h$. Alors $v\notin S_h$, et la ligne $v$ de $B_h$ est
 $e_v|_{S_h}=0\in\cR_{t,h}$.
\end{itemize}
Dans les cas (a) et (b), la condition~\eqref{eq:three} appliquée à l'occurrence $(e,v)$
donne
\[
P_v(t)|_{S_h}=P_{e,v}(t)|_{S_h}\in\cR_{t,h};
\]
c'est ici que l'uniformité par variable est utilisée. Nous raisonnons maintenant comme à
l'étape~2 de la démonstration du Lemme~\ref{lem:det}, avec $J_{t,h}$ à la place de
$I^\omega_{t,h}$. Nous obtenons $\rank B_h\le\min(a_h,r_h+|J_{t,h}|)$ et
\begin{equation}\label{eq:Dt}
\dim\ker\Phi(t)\ \ge\ \sum_h\bigl(a_h-r_h-|J_{t,h}|\bigr)_+\ \ge\ D_t-\rho_t,
\qquad D_t:=\sum_h\bigl(a_h-|J_{t,h}|\bigr)_+ .
\end{equation}

\paragraph{Un bon point.} Les étapes~1, 3 et~4 de la démonstration du Lemme~\ref{lem:det}
s'appliquent à $\Phi$. Avec~\eqref{eq:Dt} et le Lemme~\ref{lem:hermite}, elles donnent
\[
\sum_{t\in X}(D_t-\rho_t)_+\ \le\ \sum_{t\in X}\dim\ker\Phi(t)\ \le\ \nocc R\ \le\ MR
\ <\ \sum_{t\in X}(q_t+1).
\]
Il existe donc $t\in X$ avec $(D_t-\rho_t)_+<q_t+1$, c'est-à-dire $D_t<\rho_t+q_t+1$. Fixons
ce $t$.

\paragraph{Minoration de $D_t$.} Appliquons la distinction de cas du Lemme~\ref{lem:split}.

\emph{Cas (1).} $\Sigma_t$ est un $2/3$-séparateur. Comme les $J_{t,h}$ forment une partition
de $[n]\setminus\Sigma_t$,
\[
D_t\ \ge\ \sum_h\bigl(a_h-|J_{t,h}|\bigr)=n-(n-|\Sigma_t|)=|\Sigma_t|\ \ge\ \sep(F).
\]

\emph{Cas (2).} $|J_{t,h_0}|>2n/3$ pour un certain $h_0$. Alors
$\sum_{h\ne h_0}|J_{t,h}|<n/3$, et
\[
D_t=n-\sum_h\min\bigl(a_h,|J_{t,h}|\bigr)\ \ge\ n-a_{h_0}-\sum_{h\ne h_0}|J_{t,h}|
\ >\ \frac{2n}{3}-s_t .
\]
Dans les deux cas, $\rho_t+q_t+1>D_t\ge\min(\sep(F),2n/3-s_t)$, avec inégalité stricte. \qed

\subsection{Démonstration du Théorème~\ref{thm:general} (déformations générales)}

\paragraph{Passage à la moyenne dans le lemme du déterminant.} Soit $\omega$ une application
témoin uniformément aléatoire et écrivons $\E\,\dfc_t=\E\,\dfc^\omega_t$. Pour chaque $t$,
la fonction $y\mapsto(y-\rho_t)_+$ est convexe, donc par l'inégalité de Jensen
$(\E\,\dfc_t-\rho_t)_+\le\E(\dfc^\omega_t-\rho_t)_+$. En sommant sur $t$, et en utilisant le
Lemme~\ref{lem:det} pour chaque $\omega$, la Remarque~\ref{rem:mn} et le
Lemme~\ref{lem:hermite},
\[
\sum_{t\in X}\bigl(\E\,\dfc_t-\rho_t\bigr)_+\ \le\ \E\sum_{t\in X}\bigl(\dfc^\omega_t-\rho_t\bigr)_+
\ \le\ \nocc R\ \le\ MR\ <\ \sum_{t\in X}(q_t+1).
\]
Il existe donc un point $t\in X$ avec $(\E\,\dfc_t-\rho_t)_+<q_t+1$, c'est-à-dire
\begin{equation}\label{eq:goodpoint}
\E\,\dfc_t\ <\ \rho_t+q_t+1 .
\end{equation}

\paragraph{Minoration en ce point.} Appliquons le Lemme~\ref{lem:split} en ce $t$.

\emph{Cas (1).} Alors $\sum_{v\in\Sigma_t}\theta(d_v)\ge\sepw(F)$ ; en particulier
$\sepw(F)<+\infty$. La seconde borne du Lemme~\ref{lem:expdef} donne
\[
\E\,\dfc_t\ \ge\ \frac{c_0}{2\sqrt{2s_t}}\,\sepw(F)=\frac{c_0}{2\sqrt2}\cdot\frac{\sepw(F)}{\sqrt{s_t}} .
\]

\emph{Cas (2).} La première borne du Lemme~\ref{lem:expdef} et le Lemme~\ref{lem:split}(2)
donnent $\E\,\dfc_t>2n/3-s_t$.

Dans les deux cas, $\E\,\dfc_t$ est au moins égal au minimum figurant dans le théorème ;
rappelons la convention selon laquelle ce minimum vaut $2n/3-s_t$ lorsque
$\sepw(F)=+\infty$. Avec~\eqref{eq:goodpoint}, cela donne le premier énoncé.

\paragraph{Occurrences bornées.} Supposons $d_v\le b$ pour tout $v$, avec $b\ge2$.
Alors~\eqref{eq:sepw} donne $\sepw(F)\ge\theta(b)\sep(F)$, et
\[
\frac{c_0}{2\sqrt2}\cdot\frac{2(b-1)}{b^2}=c_b,\qquad
c_3=\frac{1-e^{-2}}{4}\cdot\frac{2}{9\sqrt2}=\frac{1-e^{-2}}{18\sqrt2}=0.03396\ldots
\]
Cela donne le second énoncé. \qed

\subsection{Démonstration du Lemme~\ref{lem:random} (3-CNF aléatoires)}

Nous utilisons deux propriétés de la formule aléatoire, qui sont les Lemmes~A.1 et~A.2
de~\cite{DL} ; leurs démonstrations sont reproduites à l'Annexe~\ref{app:random}. Avec
probabilité $1-o(1)$ :
\begin{itemize}
\item[(A1)] pour tout ensemble $S$ de variables avec $n/3\le|S|\le2n/3$, au moins $m/3$
 clauses contiennent à la fois une variable de $S$ et une variable hors de $S$ ;
\item[(A2)] pour tout $j\ge1$, tout ensemble de $j$ variables rencontre au plus
 $C_\Delta j\log(en/j)$ clauses, où $C_\Delta=3e^2\Delta$.
\end{itemize}
Supposons que (A1) et (A2) sont vérifiées.

\paragraph{Borne sur $\sep(F)$.} Soit $Y$ un $2/3$-séparateur avec $|Y|\le n/9$. Les
composantes de $\Gr(F)-Y$ ont chacune au plus $2n/3$ sommets et au moins $8n/9$ sommets au
total. Choisissons une union $A$ de composantes de la manière suivante. Si une composante a
au moins $n/3$ sommets, soit $A$ cette composante. Sinon, chaque composante a moins de $n/3$
sommets ; ajoutons les composantes une à une jusqu'à ce que l'union ait au moins $n/3$
sommets, et soit $A$ cette union, qui a moins de $2n/3$ sommets. Dans les deux cas,
$n/3\le|A|\le2n/3$.

Soit $e$ une clause qui contient une variable de $A$ et une variable hors de $A$. Si $e$ ne
contenait aucune variable de $Y$, ses variables seraient deux à deux adjacentes dans
$\Gr(F)-Y$, donc dans une même composante, ce qui est impossible car $A$ est une union de
composantes. Donc toute clause de ce type contient une variable de $Y$. D'après (A1), il y a
au moins $m/3>0$ telles clauses ; en particulier $Y\ne\emptyset$. D'après (A2) avec
$j=|Y|\ge1$,
\[
\frac{\Delta n}{3}=\frac m3\ \le\ C_\Delta\,|Y|\log\frac{en}{|Y|} .
\]
Avec $y=|Y|/n\in(0,1/9]$, cela s'écrit $\Delta/(3C_\Delta)\le y\log(e/y)$. La fonction
$y\mapsto y\log(e/y)$ est croissante sur $(0,1]$ et tend vers $0$ lorsque $y\to0$. Il existe
donc une constante $y^*_\Delta\in(0,1/9]$ telle que $y\log(e/y)<\Delta/(3C_\Delta)$ pour tout
$0<y<y^*_\Delta$, et par conséquent $y\ge y^*_\Delta$. Posons
$\eta_\Delta=y^*_\Delta\le1/9$. Nous avons montré que tout $2/3$-séparateur a plus de $n/9$
sommets ou au moins $\eta_\Delta n$ sommets. Donc $\sep(F)\ge\eta_\Delta n$.

\paragraph{Borne sur $\sepw(F)$.} Chaque clause a exactement trois variables distinctes,
donc $\sum_vd_v=M=3\Delta n$, et pour tout entier $D\ge1$ au plus $3\Delta n/D$ variables
vérifient $d_v>D$. Prenons $D=\lceil6\Delta/\eta_\Delta\rceil\ge2$, de sorte que
$3\Delta n/D\le\eta_\Delta n/2$. Soit $Y$ un $2/3$-séparateur avec $d_v\ge2$ pour tout
$v\in Y$. Alors $|Y|\ge\eta_\Delta n$, donc $Y$ contient au moins $\eta_\Delta n/2$
variables avec $2\le d_v\le D$, et chacune d'elles vérifie $\theta(d_v)\ge\theta(D)$ car
$\theta$ est décroissante sur $d\ge2$. Donc
$\sum_{v\in Y}\theta(d_v)\ge\frac{\eta_\Delta}{2}\theta(D)\,n$. Posons
$\eta'_\Delta=\frac{\eta_\Delta}{2}\theta(D)$. Comme $\theta\le\frac12$ sur $d\ge2$,
$\eta'_\Delta\le\eta_\Delta/4$. S'il n'existe aucun tel $Y$, $\sepw(F)=+\infty$. Dans les
deux cas, $\sepw(F)\ge\eta'_\Delta n$. \qed

\subsection{Démonstration du Corollaire~\ref{cor:main} et de la Remarque~\ref{rem:balanced}}
\label{sec:proofcor}

\paragraph{Une inégalité commune.} Dans la famille~(i), le second énoncé du
Théorème~\ref{thm:general} et $\sep(F)\ge cn$ donnent un point $t\in X$ tel que
\begin{equation}\label{eq:star}
\rho_t+q_t+1\ >\ \min\Bigl(\frac{\gamma n}{\sqrt{s_t}},\ \frac{2n}{3}-s_t\Bigr),
\end{equation}
avec $\gamma=c_bc$. Dans la famille~(ii), le premier énoncé du Théorème~\ref{thm:general} et
$\sepw(F)\ge\eta'_\Delta n$ donnent~\eqref{eq:star} avec
$\gamma=\frac{c_0}{2\sqrt2}\eta'_\Delta$. Dans la famille~(i), $c\le1$ car $\sep(F)\le n$,
et $c_b\le c_0/(4\sqrt2)<0.04$ car $(b-1)/b^2\le1/4$. Dans la famille~(ii),
$\eta'_\Delta\le\eta_\Delta\le1/9$. Donc $\gamma\le1/6$ et $c\le1$ dans les deux cas.

\paragraph{Partie 1.} Supposons $s_t\le\sqrt n\log^\ell n$ pour tout $t\in X$, et soit $t$
le point de~\eqref{eq:star}. Alors
$\gamma n/\sqrt{s_t}\ge\gamma n^{3/4}/\log^{\ell/2}n$ et
$2n/3-s_t\ge2n/3-\sqrt n\log^\ell n$. Pour $n$ grand (en fonction de $\gamma$ et de $\ell$),
la seconde quantité est au moins $n/2\ge\gamma n^{3/4}/\log^{\ell/2}n$, donc
$\rho_t+q_t+1>\gamma n^{3/4}/\log^{\ell/2}n$. Pour $n$ grand, ceci dépasse
$2\sqrt n\log^\ell n+1$, donc $\rho_t$ et $q_t$ ne peuvent pas être tous deux au plus égaux
à $\sqrt n\log^\ell n$.

\paragraph{Partie 2.} Pour $1\le j\le n/2$, on a
$\binom nj=\prod_{i=0}^{j-1}\frac{n-i}{j-i}\ge(n/j)^j\ge2^j$, et $\binom n0=2^0$. Donc, en
chaque point $t$, la charge de type produit est au moins
\begin{equation}\label{eq:charge}
\max\bigl(2^{s_t},\ 2^{\rho_t},\ 2^{\min(q_t,\lfloor n/2\rfloor)}\bigr).
\end{equation}
Soit $t$ le point de~\eqref{eq:star} et soit $n\ge216$. Si $s_t\ge n^{2/3}$, la charge est
au moins $2^{n^{2/3}}$. Sinon $s_t<n^{2/3}\le n/6$, donc $2n/3-s_t>n/2$ et
$\gamma n/\sqrt{s_t}>\gamma n^{2/3}$ ; comme $\gamma n^{2/3}\le n/2$, \eqref{eq:star} donne
$\rho_t+q_t+1>\gamma n^{2/3}$, d'où $\max(\rho_t,q_t)>(\gamma n^{2/3}-1)/2$. Si
$\rho_t>(\gamma n^{2/3}-1)/2$, la charge est au moins $2^{\rho_t}$. Si
$q_t>(\gamma n^{2/3}-1)/2$, alors, comme $(\gamma n^{2/3}-1)/2\le\lfloor n/2\rfloor$, la
charge est au moins $2^{\min(q_t,\lfloor n/2\rfloor)}\ge2^{(\gamma n^{2/3}-1)/2}$. Comme
$\gamma<1$, dans tous les cas la charge est au moins $2^{(\gamma n^{2/3}-1)/2}$.

\paragraph{Partie 3.} Le Théorème~\ref{thm:uniform} donne un point $t$ avec
$\rho_t+q_t+1>\min(\sep(F),2n/3-s_t)$, et $\sep(F)\ge cn$ dans les deux familles (d'après le
Lemme~\ref{lem:random} dans la famille~(ii)) ; la borne $b$ n'est pas utilisée. Si
$s_t\ge n/6$, la charge~\eqref{eq:charge} est au moins $2^{n/6}$. Sinon $2n/3-s_t>n/2$, donc
$\rho_t+q_t+1>c'n$ avec $c'=\min(c,1/2)$, et $\max(\rho_t,q_t)>(c'n-1)/2$. Comme
$(c'n-1)/2\le\lfloor n/2\rfloor$, \eqref{eq:charge} est au moins $2^{(c'n-1)/2}$, comme
dans la partie~2. Donc la charge est au moins $2^{\min(n/6,(c'n-1)/2)}$.

\paragraph{Remarque~\ref{rem:balanced}.} Soit $t$ le point de~\eqref{eq:star} et supposons
$s_t\le n/2$ et $s_t\le Kn/g_t$. Comme $\gamma\le1/6$ et $s_t\ge1$,
$\gamma n/\sqrt{s_t}\le n/6\le2n/3-s_t$, donc~\eqref{eq:star} s'écrit
$\rho_t+q_t+1>\gamma n/\sqrt{s_t}\ge\frac{\gamma}{\sqrt K}\sqrt{ng_t}$. Si
$q_t+1<\frac{\gamma}{2\sqrt K}\sqrt{ng_t}$, alors
$\rho_t>\frac{\gamma}{2\sqrt K}\sqrt{ng_t}$, et en divisant par $g_t$,
$\rho_t/g_t>\frac{\gamma}{2\sqrt K}\sqrt{n/g_t}\ge\frac{\gamma}{2\sqrt K}\sqrt{s_t/K}
=\frac{\gamma}{2K}\sqrt{s_t}$. \qed

\section{Le modèle de type PD est dégénéré}\label{sec:dp}

Le modèle de type PD de \cite[Section~5.2]{DL} facture $2^{s_t}$, $2^{w_t}$ et le facteur de
Taylor, mais pas $2^{\rho_t}$ (Section~\ref{sec:cost}). L'observation suivante montre que sa
charge est petite sur toute formule, pour un schéma trivial.

\begin{proposition}[Un schéma trivial dans le modèle de type PD]\label{prop:dp}
Soit $F$ une formule CNF quelconque sur $n\ge1$ variables et soit $R\ge0$. Il existe un
schéma pour $F$ au sens de la Définition~5.1 de~\cite{DL}, avec la borne $R$ sur le degré,
tel qu'en tout point $t\in X$
\[
s_t=\lceil\sqrt n\,\rceil,\qquad w_t\le\lceil\sqrt n\,\rceil-1,\qquad q_t=0,\qquad\rho_t=n .
\]
Sa charge de type PD vaut exactement $2^{\lceil\sqrt n\,\rceil}$, et sa charge de type
produit vaut $2^n$.
\end{proposition}

\begin{proof}
\emph{(a) Déformation.} Posons $P_o(T)=e_v$, un vecteur constant, pour toute occurrence
$o=(e,v)$. Alors $\deg P_o=0\le R$ et $P_o(0)=e_v$. Le polynôme $H$ de~\eqref{eq:H} est la
constante $\#F$, et $\deg H\le MR$.

\emph{(b) Reconstruction.} Soit $X=\{1,2,\dots,MR+1\}\subset\Q\setminus\{0\}$ et $q_t=0$
pour tout $t\in X$. Les données de Hermite de $f$ sont les valeurs $(f(t))_{t\in X}$. Soit
$\Lambda\bigl((y_t)_{t\in X}\bigr)=\sum_{t\in X}y_t\,L_t(0)$, où
$L_t(T)=\prod_{u\in X\setminus\{t\}}\frac{T-u}{t-u}$ est la base de Lagrange. C'est une
fonctionnelle linéaire, et $\Lambda(f)=f(0)$ pour tout polynôme $f$ de degré au plus
$|X|-1=MR$, par interpolation de Lagrange.

\emph{(c) Designs.} Soit $s=\lceil\sqrt n\,\rceil$ et $g=\lceil n/s\rceil$. En tout $t\in X$,
utilisons le même design : les groupes $S_1,\dots,S_g$ sont des blocs consécutifs de $[n]$,
les $g-1$ premiers de taille $s$ et le dernier de taille au plus $s$ ; $\alpha_t(e)$ est un
groupe arbitraire ; et $\cR_{t,h}=\Q^{S_h}$ pour tout $h$. La condition~\eqref{eq:three} est
vérifiée parce que tout vecteur de $\Q^{S_h}$ appartient à $\cR_{t,h}$.

\emph{Paramètres.} Comme $1\le s\le n$, le premier groupe a exactement $s$ éléments, donc
$s_t=s$. Comme $n/s\le n/\sqrt n=\sqrt n\le s$ et que $s$ est un entier, $g\le s$. Le graphe
d'interaction $\cI_t$ a $g$ sommets, donc $w_t=\tw(\cI_t)\le g-1\le s-1$. Enfin
$\rho_t=\sum_h|S_h|=n$ et $\binom{n}{\min(q_t,\lfloor n/2\rfloor)}=\binom n0=1$. La charge de
type PD est $\max(2^{s},2^{w_t},1)=2^{s}$, et la charge de type produit est
$\max(2^s,2^{n})=2^n$.
\end{proof}

\paragraph{Conséquences.}
\begin{enumerate}
\item \emph{La borne de~\cite{DL} est optimale dans le modèle de type PD, pour une raison
 triviale.} Pour toute formule, la plus petite charge de type PD d'un schéma est au plus
 $2^{\lceil\sqrt n\,\rceil}$. Avec \cite[Théorème~8.1]{DL}, sur toute famille de largeur
 arborescente primale au moins $cn$, elle vaut $2^{\Theta(\sqrt n)}$ : pour $n$ grand,
 l'exposant est compris entre $\sqrt{(cn+1)/2}-1$ et $\lceil\sqrt n\,\rceil$. Les théorèmes
 de~\cite{DL} sont des bornes inférieures et restent corrects. Ce que montre la
 Proposition~\ref{prop:dp}, c'est qu'aucun argument, par quelque technique que ce soit, ne
 peut démontrer plus que $2^{\lceil\sqrt n\,\rceil}$ dans ce modèle. La phrase de
 \cite[Section~9]{DL} selon laquelle \eng{``our lower-bound technique cannot exceed
 $2^{\tilde O(\sqrt n)}$''}\footnote{Traduction : \og notre technique de borne inférieure
 ne peut pas dépasser $2^{\tilde O(\sqrt n)}$\fg{}.} exprime donc une propriété du modèle et
 pas seulement de la technique.
\item \emph{L'écart de \cite[Section~9]{DL} n'est pas ouvert dans le modèle de type PD.} La
 cible de coût $2^{\tilde O(\sqrt n)}$ \eng{``from above''}\footnote{Traduction : \og par
 le haut\fg{}.} y est atteinte par le schéma trivial, avec $\sqrt n$ lignes par groupe au
 lieu d'un nombre polylogarithmique. La question de savoir si les mesures peuvent absorber
 la coupe n'a de sens que dans un modèle qui facture les lignes. C'est le cas du modèle de
 type produit, et le Corollaire~\ref{cor:main} s'y applique ; le schéma trivial a une
 charge de type produit de $2^n$, en accord avec le Théorème~\ref{thm:uniform}.
\item \emph{Pourquoi le modèle sous-facture.} La charge $2^{w_t}$ compte deux valeurs par
 groupe mesuré, quel que soit son nombre de lignes ; c'est la justification citée à la
 Section~\ref{sec:cost}. Dans le schéma trivial, un groupe est entièrement mesuré, et sa
 lecture détermine toute son affectation, donc elle prend $2^{|S_h|}$ valeurs. Une
 programmation dynamique sur les lectures serait alors une programmation dynamique sur les
 affectations de blocs de variables, ce que la charge $2^{w_t}$ ne tarifie pas. Un modèle
 de type PD réparé doit pondérer un groupe par son nombre de lignes ; voir la
 Section~\ref{sec:concl}.
\end{enumerate}

\section{Vérification par machine}\label{sec:lean}

La partie déterministe de l'argument est formalisée en Lean~4~\cite{Lean4} avec
Mathlib~\cite{Mathlib}. Les sources se trouvent dans le répertoire \texttt{lean/} du dépôt
\begin{center}
\url{https://github.com/luluddof/pnp-recherche}
\end{center}
La chaîne d'outils est Lean \texttt{v4.37.0-rc1}, et le commit de Mathlib est
\texttt{e62b985f}, tels que fixés dans le fichier \texttt{lake-manifest.json}. Les quatre
fichiers sources \texttt{WitnessDet}, \texttt{Scheme}, \texttt{TheoremB} et
\texttt{TheoremA} (dans \texttt{lean/PnpLean/}) ne contiennent aucun \texttt{sorry} et ne
déclarent aucun axiome. Pour reproduire :
\begin{verbatim}
  cd lean && lake exe cache get && lake build PnpLean.TheoremA
\end{verbatim}
La cible \texttt{PnpLean.TheoremA} importe les trois autres fichiers. Nous avons exécuté la
commande de compilation avec ces versions (le cache de Mathlib étant déjà présent) et elle
s'est terminée avec succès. Nous avons ensuite exécuté la commande \texttt{\#print axioms}
sur dix déclarations du Tableau~\ref{tab:lean} :
\begin{center}
\texttt{sum\_nullity\_le}, \texttt{deficiency\_le\_nullity\_add}, \texttt{lemma1},
\texttt{factR}, \texttt{avg\_lemma1},\\ \texttt{exists\_point}, \texttt{lemma3},
\texttt{theoremB}, \texttt{theoremB\_bounded}, \texttt{theoremA}.
\end{center}
Pour chacune de ces déclarations, la commande a signalé exactement les trois axiomes
standard \texttt{propext}, \texttt{Classical.choice} et \texttt{Quot.sound}.

\begin{table}[tb]
\centering
\small
\begin{tabular}{@{}>{\raggedright\arraybackslash}p{3.6cm}>{\raggedright\arraybackslash}p{6.1cm}>{\raggedright\arraybackslash}p{4.9cm}@{}}
\toprule
Énoncé de cet article & Déclaration Lean & Statut\\
\midrule
Lemme~\ref{lem:hermite} (comptage de Hermite) & \texttt{Scheme.factR} & formalisé\\[2pt]
Lemme~\ref{lem:det}, étape 1 & \texttt{det\_ne\_zero\_of\_eval\_zero}, \newline\texttt{natDegree\_det\_le} & formalisé avec la borne $nR$ (et non $\nocc R$)\\[2pt]
Lemme~\ref{lem:det}, étape 2 & \texttt{deficiency\_le\_nullity\_add}, \newline\texttt{Scheme.deficiency\_sub\_le\_nullity} & formalisé\\[2pt]
Lemme~\ref{lem:det}, étapes 3--4 & \texttt{pow\_dvd\_det\_of\_nullity}, \newline\texttt{sum\_nullity\_le} & formalisé\\[2pt]
Lemme~\ref{lem:det} & \texttt{Scheme.lemma1} & formalisé avec la borne $nR$\\[2pt]
Étape de passage à la moyenne du Théorème~\ref{thm:general}, et~\eqref{eq:goodpoint} & \texttt{Scheme.avg\_lemma1}, \newline\texttt{Scheme.exists\_point} & formalisé ; \texttt{exists\_point} suppose $n\le M$\\[2pt]
Lemme~\ref{lem:split} & \texttt{Scheme.lemma3} & formalisé, sous la forme de la disjonction \og $\Sigma_t$ est un $2/3$-séparateur, ou $\frac12\sum_h|\Nbar_{t,h}-a_h|>2n/3-s_t$\fg{}\\[2pt]
Théorème~\ref{thm:uniform} & \texttt{Scheme.theoremA} & formalisé, sans condition, en supposant $n\le M$\\[2pt]
Théorème~\ref{thm:general} & \texttt{Scheme.theoremB}, \newline\texttt{Scheme.theoremB\_bounded} & formalisé \emph{conditionnellement} : le Lemme~\ref{lem:expdef} est une hypothèse explicite \texttt{Lemma2Hyp} ; suppose $n\le M$\\
\midrule
Lemme~\ref{lem:anti}, Lemme~\ref{lem:expdef} & --- & \textbf{non formalisé} (démonstration sur papier uniquement)\\[2pt]
Moyenne uniforme $=$ choix uniformes indépendants & --- & \textbf{non formalisé} (argumenté sur papier, Section~\ref{sec:ournot})\\[2pt]
Borne $\nocc R\le MR$ sans $n\le M$ (Remarque~\ref{rem:mn}) & --- & \textbf{non formalisé}\\[2pt]
Lemme~\ref{lem:random}, Corollaire~\ref{cor:main}, Remarque~\ref{rem:balanced}, Proposition~\ref{prop:dp}, Annexe~\ref{app:random} & --- & \textbf{non formalisé}\\
\bottomrule
\end{tabular}
\caption{Ce qui est vérifié par machine et ce qui ne l'est pas. Les déclarations sans le
préfixe \texttt{Scheme.} sont dans l'espace de noms \texttt{WitnessDet} ; les autres sont
dans \texttt{ClauseScheme.Scheme}. Les noms Lean \texttt{lemma1}, \texttt{lemma3},
\texttt{theoremA}, \texttt{theoremB}, \texttt{factR} suivent une numérotation antérieure des
résultats.}\label{tab:lean}
\end{table}


\paragraph{Ce que disent les théorèmes Lean.} Dans \texttt{Scheme.theoremB}, l'espérance
$\E\,\dfc_t$ est définie comme la moyenne arithmétique de $\dfc^\omega_t$ sur le type fini
de toutes les applications témoins, et \texttt{Lemma2Hyp c0} est l'énoncé du
Lemme~\ref{lem:expdef} pour cette moyenne, avec une constante $c_0\ge0$ laissée en
paramètre. Le théorème affirme : si $n\le M$, si \texttt{Lemma2Hyp c0} est vraie, et si $W$
est un nombre réel tel que $W\le\sum_{v\in Y}\theta(d_v)$ pour tout $2/3$-séparateur $Y$
dont toutes les variables vérifient $d_v\ge2$, alors un certain $t\in X$ vérifie
$\min\bigl(\frac{c_0}{2\sqrt2}W/\sqrt{s_t},\,2n/3-s_t\bigr)<\rho_t+q_t+1$. La vérification
par machine couvre donc le squelette logique du Théorème~\ref{thm:general} (déterminant,
comptage de Hermite, passage à la moyenne, lemme du séparateur, analyse de cas et
constantes) et \emph{non} son ingrédient probabiliste. \texttt{Scheme.theoremA} énonce le
Théorème~\ref{thm:uniform} sans autre hypothèse non démontrée que $n\le M$.

\paragraph{Différences entre les définitions Lean et celles de l'article.}
\begin{enumerate}
\item \emph{Formule.} Une formule est donnée uniquement par les ensembles de variables de ses
 clauses $I:\kappa\to\mathrm{Finset}(\mathrm{Fin}\ n)$ pour un type fini $\kappa$ de
 clauses. Les signes, les polynômes $\varphi_e$ et $H$ n'apparaissent pas. La structure
 \texttt{Scheme} enregistre ce que les démonstrations utilisent : les vecteurs $P_o$ à
 coefficients de degré au plus $R$ avec $P_o(0)=e_v$ ; un ensemble $X$ de rationnels ne
 contenant pas $0$ ; des ordres $q$ ; et une fonctionnelle linéaire $\Lambda$ sur les
 données de Hermite, exacte sur tout polynôme de degré au plus $MR$. C'est un ensemble
 d'hypothèses plus faible que la Définition~5.1, donc les théorèmes Lean s'appliquent à tout
 schéma de l'article.
\item \emph{Groupes.} Les groupes sont indexés par un même type fini fixé $G$ en tous les
 points, et la partition en $t$ est une application des variables vers $G$ ; les groupes
 vides sont autorisés. Un schéma de l'article est codé en prenant $G$ assez grand et en
 attribuant aux étiquettes inutilisées le groupe vide et le sous-espace nul. Les données de
 design sont des fonctions de tout $t\in\Q$, mais la condition~\eqref{eq:three} n'est
 exigée que pour $t\in X$.
\item \emph{Sous-espaces de mesure.} $\cR_{t,h}$ est un sous-espace de $\Q^n$ dont les
 vecteurs s'annulent hors de $S_{t,h}$, au lieu d'un sous-espace de $\Q^{S_{t,h}}$ ; la
 condition~\eqref{eq:three} est énoncée pour le vecteur $P_{e,v}(t)$ dont les coordonnées
 hors de $S_{t,h}$ sont mises à zéro.
\item \emph{Séparateurs.} \texttt{IsSep I Y} dit que pour tout $v\notin Y$, l'ensemble des
 sommets atteignables depuis $v$ dans le graphe primal privé de toutes les arêtes touchant
 $Y$ a au plus $2n/3$ éléments. Ceci est équivalent à la Définition~\ref{def:sep}. Les
 nombres $\sep(F)$ et $\sepw(F)$ ne sont pas définis ; les théorèmes prennent une borne
 inférieure arbitraire (\texttt{sepLB}, respectivement $W$) valable pour tout séparateur, ce
 qui donne les énoncés de l'article en prenant la meilleure borne inférieure.
\item \emph{Hypothèse $n\le M$.} En Lean, les énoncés de \texttt{exists\_point}, de
 \texttt{theoremA}, de \texttt{theoremB} et de \texttt{theoremB\_bounded} supposent
 $n\le M$ et utilisent $\deg\det\Phi^\omega\le nR\le MR$. L'article supprime cette hypothèse grâce à la
 borne plus fine $\deg\det\Phi^\omega\le\nocc R$ (étape~1 de la démonstration du
 Lemme~\ref{lem:det}) ; ce raffinement d'une ligne n'est pas formalisé.
\item \emph{Arithmétique.} Les déficits et les parties positives sont des soustractions
 tronquées d'entiers naturels ; $s_t$ est le supremum de $a_h$ sur $G$ ; dans
 \texttt{theoremB\_bounded}, la borne sur les occurrences vérifie $b\ge2$.
\end{enumerate}
La correspondance entre les énoncés Lean et les énoncés de l'article a été vérifiée par
lecture, par des agents d'IA et non par un expert humain ; un décalage entre une définition
Lean et sa signification voulue est le principal risque résiduel de cette section.

\section{Vérifications par le calcul}\label{sec:checks}

Nous avons exécuté quatre petits scripts. Ils testent l'algèbre et les inégalités sur de
petites instances. \textbf{De telles vérifications ne peuvent que détecter des erreurs ;
elles ne démontrent rien, et en particulier rien au sujet des énoncés asymptotiques.} Les
scripts se trouvent dans le dépôt cité à la Section~\ref{sec:lean}, dans
\texttt{code/gen6/g6a/} et \texttt{code/gen6/g6r/}. Nous ne rapportons que ce que nous avons
observé en réexécutant des copies non modifiées de ces scripts pour cette version, avec
Python~3.13.18, sympy~1.14.0 et networkx~3.7, sur une machine Linux (un seul fil
d'exécution). Le Tableau~\ref{tab:checks} donne les commandes, exécutées depuis la racine du
dépôt.

\begin{table}[htbp]
\centering
\small
\begin{tabular}{@{}llr@{}}
\toprule
Vérification & Commande & Temps\\
\midrule
C1 & \texttt{cd code/gen6/g6a \&\& python3 run\_checks.py theorem} & 6 s\\
C2 & \texttt{cd code/gen6/g6a \&\& python3 run\_checks.py e2e} & 19 s\\
C3 & \texttt{cd code/gen6/g6a \&\& python3 anticonc\_check.py} & 0.5 s\\
C4 & \texttt{cd code/gen6/g6r \&\& python3 lemma1\_check.py} & 125 s\\
\bottomrule
\end{tabular}
\caption{Commandes des vérifications par le calcul. C1 nécessite networkx ; C2 et C4
nécessitent sympy. Tous les scripts utilisent des graines fixées.}\label{tab:checks}
\end{table}

\paragraph{C1 : le Lemme~\ref{lem:det} et la borne du Lemme~\ref{lem:expdef}.}
\emph{Instances.} Des formules aléatoires à $n$ variables et $n$ clauses dans lesquelles
chaque variable apparaît dans exactement trois clauses, pour $(n,|X|,R)=(24,8,2)$,
$(30,10,2)$ et $(36,9,3)$, avec $X=\{1,\dots,|X|\}$ et des groupes de taille $6$ dans une
partition aléatoire en chaque point. Quatre familles de designs : couper-et-mesurer avec une
déformation uniforme par variable ; couper-et-mesurer avec une déformation dépendant de la
clause ; un design \og à somme globale\fg{} ; une déformation dense aléatoire avec le plus
petit $\cR_{t,h}$ admissible. Cela fait $12$ exécutions.

\emph{Ce qui est testé par assertion.} La condition~\eqref{eq:three}, exactement sur $\Q$.
Pour quatre applications témoins aléatoires par exécution :
$\dim\ker\Phi^\omega(t)\ge\dfc^\omega_t-\rho_t$ en chaque point, et
$\sum_t\dim\ker\Phi^\omega(t)\le nR$, avec des rangs exacts sur $\Q$. En chaque point, que
$\E\,\dfc_t$ est au moins égal à la borne du Lemme~\ref{lem:expdef} avec $c_0$ remplacé par
$0.216$ ; ici $\E\,\dfc_t$ est calculé à partir de la distribution exacte de chaque $N_h$ en
arithmétique à virgule flottante. Enfin $\sum_t(\E\,\dfc_t-\rho_t)_+\le nR$.

\emph{Ce que nous avons observé.} Toutes les assertions ont été satisfaites dans les $12$
exécutions. Dans l'ordre des quatre familles, la plus grande valeur de
$\sum_t\dim\ker\Phi^\omega(t)$ sur les quatre applications témoins était $48$, $48$, $40$,
$0$ pour $n=24$ (où $nR=48$) ; $60$, $60$, $50$, $0$ pour $n=30$ (où $nR=60$) ; et $108$,
$108$, $99$, $0$ pour $n=36$ (où $nR=108$). La borne $nR$ du Lemme~\ref{lem:det} a donc été
atteinte par les designs couper-et-mesurer. Sur tous les points de toutes les exécutions,
$\E\,\dfc_t$ variait de $2.57$ à $8.46$ et $\rho_t$ de $10$ à $36$.

\emph{Limites de cette vérification.} Dans chacune des $12$ exécutions, $\rho_t$ dépassait
$\E\,\dfc_t$ en chaque point (d'au moins $6.9$), donc $\sum_t(\E\,\dfc_t-\rho_t)_+=0$.
\textbf{L'étape de passage à la moyenne du Théorème~\ref{thm:general} n'est donc jamais
réellement mise à l'épreuve} : dans ces exécutions, l'inégalité
$\sum_t(\E\,\dfc_t-\rho_t)_+\le nR$ est trivialement vraie. Ces designs ont $\rho_t$ compris entre $0.4n$ et $n$
et sont loin du régime du Théorème~\ref{thm:general} ; ils ne testent ni l'optimalité du
facteur $\sqrt s$ ni les inégalités finales des Théorèmes~\ref{thm:uniform}
et~\ref{thm:general}. Un fichier de sortie conservé dans le dépôt
(\texttt{theorem\_checks.txt}) ne contient que les configurations $n=24$ et $n=30$ ; dans
notre exécution, la configuration $n=36$ s'est terminée normalement, et la sortie pour
$n=24$ et $n=30$ était identique au fichier conservé.

\paragraph{C2 : reconstruction de bout en bout.} Une formule avec $n=12$ (trois occurrences
par variable), trois groupes de quatre variables, $R=1$, $X=\{1,\dots,6\}$, $q_t=6$ (donc
$\sum_t(q_t+1)=42>MR=36$), et deux familles de designs. Le script calcule $H(T)$ par force
brute modulo le nombre premier $2^{61}-1$, reconstruit $H(0)$ à partir des données de
Hermite, et en chaque point compare une évaluation de type produit (énumération des uplets
de lectures de mesure, comptage à l'intérieur des groupes) avec la valeur de $H(t)$ obtenue
par force brute en arithmétique rationnelle exacte.
\emph{Observé :} $\#F=977$ ; $H(0)=977$ et la valeur reconstruite par Hermite est $977$ pour
les deux designs ; l'évaluation de type produit est égale à $H(t)$ aux $6$ points pour les
deux designs, avec entre $64$ et $3072$ uplets de lectures. Ceci vérifie le code et le sens
de la Définition~5.1 sur une instance ; cela ne teste pas nos théorèmes.

\paragraph{C3 : anticoncentration.} $4000$ sommes aléatoires d'au plus $40$ variables de
Bernoulli indépendantes ; pour chacune, le minimum sur les entiers $a$ de $\E|B-a|$ divisé
par $\min(V,\sqrt V)$, en arithmétique à virgule flottante. \emph{Observé :} le plus petit
rapport est $0.7397$, contre $c_0\approx0.2162$ dans le Lemme~\ref{lem:anti}. C'est un test
de cohérence sur des instances aléatoires, et non une démonstration ni une recherche de
pires cas.

\paragraph{C4 : un test indépendant du Lemme~\ref{lem:det}.} Le script construit des schémas
aléatoires avec $n=9$, $7$ clauses de trois variables, $R=2$, $4$ points, $3$ groupes et une
ligne de mesure aléatoire par groupe, en résolvant les contraintes
linéaires~\eqref{eq:three} pour les coefficients de la déformation (sympy, arithmétique
exacte). Pour chaque schéma et six applications témoins aléatoires, il teste par assertion
$P_o(0)=e_v$, $\det\Phi^\omega(0)=1$, $\deg\det\Phi^\omega\le nR$,
$\dim\ker\Phi^\omega(t)\ge\dfc^\omega_t-\rho_t$, que la multiplicité de la racine $t$ de
$\det\Phi^\omega$ est au moins $\dim\ker\Phi^\omega(t)$, et
$\sum_t(\dfc^\omega_t-\rho_t)_+\le nR$. \emph{Observé :} parmi les six graines $0,\dots,5$,
trois (les graines $0$, $4$, $5$) n'ont produit aucun schéma, parce que le système linéaire
était incompatible, et le script n'a rien renvoyé pour elles. Pour les graines $1$, $2$,
$3$, toutes les assertions ont été satisfaites, sur $3\times6=18$ applications témoins au
total. Dans ces trois schémas, $\rho_t=3$ en chaque point, et la plus grande valeur de
$\sum_t(\dfc^\omega_t-\rho_t)_+$ était $1$, $3$ et $1$, contre $nR=18$. Le test est donc
faible : il est loin du cas d'égalité.

\paragraph{Ce qui n'est pas vérifié numériquement.} Le Lemme~\ref{lem:split}, le
Lemme~\ref{lem:random}, le Corollaire~\ref{cor:main} et la Proposition~\ref{prop:dp} n'ont
pas été vérifiés numériquement pour cette version, et aucune exécution n'exhibe un point
avec $\E\,\dfc_t>\rho_t$.

\section{Méthodologie et usage des outils d'IA}\label{sec:ai}

Nous indiquons sans détour comment ce travail a été produit.

\paragraph{Mathématiques produites par IA.} Les résultats de cet article ont été trouvés,
démontrés, rédigés et relus par des agents d'IA, des instances de Claude (Anthropic),
travaillant sous la direction de l'auteur, qui a supervisé le processus de travail et décidé
de ce qui serait publié. L'auteur n'a pas écrit les démonstrations.

\paragraph{Relectures, et leurs limites.} Un premier agent a produit l'argument du
déterminant, des démonstrations et des scripts de test. Un deuxième agent a effectué une
relecture contradictoire de chaque lemme et a écrit le test indépendant C4. Un troisième
agent a rédigé un rapport d'évaluation sur une version antérieure de ce texte ; ce rapport a
jugé les démonstrations correctes dans la limite de ce qu'il pouvait vérifier, et a exigé
les changements de portée qui figurent maintenant dans l'article : le modèle de type PD est
dégénéré plutôt qu'\og ouvert\fg{} (Section~\ref{sec:dp}), la comparaison se fait avec
\cite[Théorème~8.1 et Lemme~8.2]{DL} et non avec leur Théorème~8.3, et la section sur les
vérifications par le calcul a été corrigée. Cette version a ensuite été rédigée par un autre
agent, qui a refait les démonstrations et réexécuté les vérifications.
\textbf{Ces relectures ne sont pas indépendantes des auteurs des démonstrations} : tous les
agents sont des instances de la même famille de modèles et peuvent partager des angles
morts. Elles ne remplacent pas une relecture par des experts humains.

\paragraph{Vérification par machine.} Pour réduire ce risque, la partie déterministe de
l'argument a été formalisée en Lean~4 (Section~\ref{sec:lean}). Les démonstrations formelles
ont elles aussi été écrites par des agents d'IA, mais elles sont vérifiées par le noyau de
Lean. La partie probabiliste (Lemmes~\ref{lem:anti} et~\ref{lem:expdef}), le lemme sur les
3-CNF aléatoires, le corollaire et la Proposition~\ref{prop:dp} ne sont pas formalisés.

\paragraph{Aucune relecture par un expert humain.} Au moment de la rédaction, \textbf{aucun
expert humain n'a relu ce travail}. L'auteur est étudiant en informatique. Les lecteurs
doivent considérer les résultats en conséquence, et les corrections sont les bienvenues à
l'adresse indiquée en première page.

\paragraph{Sources et nouveauté.} Les définitions et la question sont citées d'après le
texte de~\cite{DL}. Parmi les travaux cités, les agents ont lu \cite{DL} en entier ; les
prépublications~\cite{OpenAI,Claude} et l'exposé~\cite{ODonnell} ne sont cités que pour le
contexte, et le présent article ne s'appuie pas sur leur contenu au-delà de ce que
rapporte~\cite{DL}. Les références classiques ont été fournies par les agents d'IA à partir
de leur connaissance de la littérature et n'ont pas été relues pour cet article.
L'appréciation de la nouveauté repose sur des recherches effectuées par des agents d'IA, et
non sur un état de l'art systématique.

\paragraph{Responsabilité.} L'auteur assume la responsabilité du contenu de cet article, y
compris de toute erreur qui subsisterait.

\section{Conclusion et problèmes ouverts}\label{sec:concl}

Dias et Lima ont demandé une démonstration du fait que les mesures ne peuvent pas absorber
la coupe dans les schémas d'interpolation par produit de clauses. Nous avons démontré une
forme quantitative de cet énoncé : en un certain point d'évaluation, le nombre de lignes de
mesure plus l'ordre de Taylor dépasse $c_b\sep(F)/\sqrt{s_t}$, ou bien les groupes sont trop
grands ; et pour les déformations uniformes par variable, le facteur $\sqrt{s_t}$ disparaît.
Sur les formules à occurrences bornées à séparateurs linéaires, et sur les 3-CNF aléatoires
de densité supérieure à $3$, le design décrit dans \cite[Section~9]{DL} n'existe pas, et
tout schéma appliqué à la formule elle-même a une charge de type produit de
$2^{\Omega(n^{2/3})}$. Nous avons aussi observé que le modèle de type PD ne peut certifier
plus que $2^{\lceil\sqrt n\rceil}$ sur aucune formule.

Ce sont des résultats négatifs modestes sur une abstraction d'une voie d'une prépublication.
Ils ne concernent ni la forme de parité ni la contraction de tenseurs de la prépublication,
ce sont des bornes inférieures sur les charges d'un modèle et non sur le temps d'exécution,
et ils n'ont aucune incidence sur ETH ni sur P contre NP.

\paragraph{Problèmes ouverts.}
\begin{enumerate}
\item \emph{Optimalité du facteur $\sqrt s$.} Existe-t-il une formule avec
 $\sep(F)=\Omega(n)$ et un schéma avec $\rho_t+q_t=O(n/\sqrt{s_t})$ en chaque point et
 $s_t=o(n)$ ? D'après le Théorème~\ref{thm:uniform}, sa déformation doit dépendre de la
 clause. Le facteur provient uniquement des fluctuations d'une application témoin aléatoire,
 et un autre choix de témoins, ou un autre mineur, pourrait le supprimer. S'il peut être
 supprimé, la borne de type produit devient $2^{\Omega(n)}$ ; s'il ne le peut pas, la
 vérité pourrait être aussi basse que $2^{\tilde\Theta(n^{2/3})}$.
 Nous ne connaissons aucune borne supérieure meilleure que celle de la prépublication dans
 le modèle de type produit.
\item \emph{Les entrées réelles de la prépublication.} Démontrer
 $\sep(F^\star)=\Omega(n^\star)$, ou son analogue pondéré, pour la formule à occurrences
 bornées $F^\star$ construite avec les cycles de Tovey~\cite{Tov84}, et pour les feuilles de
 la sparsification. Cela étendrait le Corollaire~\ref{cor:main} au cadre de
 \cite[Théorème~8.3]{DL}. Nous ne l'avons pas fait.
\item \emph{Un modèle de type PD pertinent.} D'après la Proposition~\ref{prop:dp}, la charge
 de type PD doit être réparée avant que la question puisse y être posée. Un candidat naturel
 remplace $2^{w_t}$ par $2^{W_t}$, où $W_t$ est le minimum, sur les décompositions
 arborescentes du graphe sur les groupes mesurés utilisé dans~\cite{DL}, du plus grand
 nombre total de lignes dans un sac. A-t-on $\max(s_t,W_t,q_t)\ge n^{1/2+\delta}$ en un
 certain point, sur les formules à séparateurs linéaires ? Cette question est ouverte. Des
 résultats partiels dans cette direction existent dans le dépôt du projet ; ils sont
 conditionnels et n'ont pas été relus, et nous ne les revendiquons pas ici.
\item \emph{Formalisation.} Formaliser les Lemmes~\ref{lem:anti} et~\ref{lem:expdef}, de
 sorte que le Théorème~\ref{thm:general} soit vérifié par machine sans hypothèse, et
 supprimer l'hypothèse $n\le M$ des énoncés Lean.
\item \emph{Constantes.} Améliorer $c_0$ ; les expériences suggèrent que le
 Lemme~\ref{lem:anti} est loin d'être optimal.
\item \emph{Au-delà de la Définition~5.1.} Des limites similaires valent-elles pour la forme
 de parité ou la contraction de tenseurs de la prépublication ? Comme expliqué dans
 \cite[Section~5.1]{DL}, celles-ci rompent le lien entre le nombre d'occurrences et le degré
 de reconstruction exacte, sur lequel reposent le Lemme~\ref{lem:hermite} et notre
 comparaison des degrés. Notre méthode ne s'y applique pas en l'état.
\end{enumerate}

\appendix
\section{Deux propriétés des 3-CNF aléatoires (d'après Dias et Lima)}\label{app:random}

Par souci de complétude, nous reproduisons les démonstrations des deux faits utilisés dans
la démonstration du Lemme~\ref{lem:random}. Ce sont les Lemmes~A.1 et~A.2 de~\cite{DL}, et
les démonstrations suivent les leurs. $F$ a $n$ variables et $m=\Delta n$ clauses
indépendantes, chacune portant sur trois variables distinctes uniformément aléatoires avec
des signes aléatoires ; $\log$ est le logarithme naturel, et \og avec forte probabilité\fg{}
signifie avec une probabilité tendant vers $1$ lorsque $n\to\infty$.

\begin{lemma}[{\cite[Lemme~A.1]{DL}}]\label{lem:A1}
Si $\Delta>3$, alors avec forte probabilité, pour tout ensemble $S$ de variables avec
$n/3\le|S|\le2n/3$, au moins $m/3$ clauses contiennent à la fois une variable de $S$ et une
variable hors de $S$.
\end{lemma}

\begin{proof}
Fixons $S$ avec $|S|=\alpha n$, $\alpha\in[1/3,2/3]$. L'ensemble des variables d'une clause
est une partie à $3$ éléments uniforme de $[n]$, donc il est inclus dans $S$ avec
probabilité $\binom{|S|}{3}/\binom n3\le\alpha^3$, et inclus dans le complémentaire de $S$
avec probabilité au plus $(1-\alpha)^3$. La fonction $\alpha^3+(1-\alpha)^3$ est convexe,
donc sur $[1/3,2/3]$ elle est au plus égale à sa valeur $1/27+8/27=1/3$ aux extrémités.
Donc chaque clause, indépendamment, contient une variable de $S$ et une variable hors de $S$
avec probabilité au moins $2/3$, et le nombre de telles clauses domine stochastiquement une
variable binomiale $\mathrm{Bin}(m,2/3)$. D'après la borne de
Chernoff--Hoeffding~\cite[Théorème~1]{Hoe63},
\[
\PP\bigl(\mathrm{Bin}(m,2/3)\le m/3\bigr)\ \le\ \exp\bigl(-m\,\mathrm{KL}(1/3\,\|\,2/3)\bigr)
=2^{-m/3},
\]
car $\mathrm{KL}(1/3\,\|\,2/3)=\frac13\ln\frac12+\frac23\ln2=\frac13\ln2$. Il y a au plus
$2^n$ ensembles $S$, donc la probabilité qu'un certain $S$ soit en défaut est au plus
$2^{n-\Delta n/3}$, qui tend vers $0$ car $\Delta>3$.
\end{proof}

\begin{lemma}[{\cite[Lemme~A.2]{DL}}]\label{lem:A2}
Soit $C_\Delta=3e^2\Delta$. Avec forte probabilité, pour tout $1\le j\le n$, tout ensemble
de $j$ variables rencontre au plus $C_\Delta\,j\log(en/j)$ clauses.
\end{lemma}

\begin{proof}
Fixons un ensemble $S$ de $j$ variables et soit $L_S$ le nombre de couples (clause, variable
de $S$ dans cette clause). Le nombre de clauses qui rencontrent $S$ est au plus $L_S$. Pour
une clause, le nombre de ses variables dans $S$ est hypergéométrique : trois tirages sans
remise parmi $n$ variables, dont $j$ sont dans $S$. D'après un théorème de
Hoeffding~\cite[Théorème~4]{Hoe63}, sa fonction génératrice des moments est au plus celle de
$\mathrm{Bin}(3,j/n)$. Les clauses sont indépendantes, donc pour tout $\lambda>0$,
\[
\E e^{\lambda L_S}\ \le\ \bigl(1+\tfrac jn(e^\lambda-1)\bigr)^{3m}\ \le\ \exp\bigl(\mu(e^\lambda-1)\bigr),
\qquad \mu=\frac{3mj}{n}=3\Delta j .
\]
Par l'inégalité de Markov avec $\lambda=\ln(x/\mu)$, pour tout $x>\mu$,
\[
\PP(L_S\ge x)\ \le\ \exp\bigl(-x\ln(x/\mu)+x-\mu\bigr)\ \le\ \Bigl(\frac{e\mu}{x}\Bigr)^{x}.
\]
Prenons $x=C_\Delta\,j\log(en/j)$. Comme $j\le n$, $\log(en/j)\ge1$, donc
$x\ge C_\Delta j>\mu$ et $e\mu/x=1/(e\log(en/j))\le1/e$. Donc
$\PP(L_S\ge x)\le e^{-x}=(en/j)^{-C_\Delta j}$. Il y a $\binom nj\le(en/j)^j$ ensembles de
taille $j$, donc la probabilité qu'un certain ensemble de taille $j$ soit en défaut est au
plus $(en/j)^{-(C_\Delta-1)j}\le(en/j)^{-2j}$. Pour $j\le\sqrt n$, ceci est au plus
$(e^2n)^{-j}$, et pour $j>\sqrt n$, c'est au plus $e^{-2j}$. En sommant sur $1\le j\le n$,
on obtient au plus $2/(e^2n)+2e^{-2\sqrt n}$, qui tend vers $0$.
\end{proof}

\begin{thebibliography}{99}
\englishhyph

\bibitem{AVW18} J. Alman, V. Vassilevska Williams. Limits on all known (and some unknown)
 approaches to matrix multiplication. \emph{FOCS} 2018.

\bibitem{AFL15} A. Ambainis, Y. Filmus, F. Le Gall. Fast matrix multiplication: limitations
 of the Coppersmith--Winograd method. \emph{STOC} 2015.

\bibitem{BSI10} E. Ben-Sasson, R. Impagliazzo. Random CNF's are hard for the polynomial
 calculus. \emph{Comput. Complexity} 19, 2010.

\bibitem{BSW01} E. Ben-Sasson, A. Wigderson. Short proofs are narrow---resolution made
 simple. \emph{J. ACM} 48(2), 2001.

\bibitem{Ber97} B. Berger. The fourth moment method. \emph{SIAM J. Comput.} 26(4), 1997.

\bibitem{Bod93} H. L. Bodlaender. A tourist guide through treewidth. \emph{Acta Cybernetica}
 11, 1993.

\bibitem{CS88} V. Chv\'atal, E. Szemer\'edi. Many hard examples for resolution.
 \emph{J. ACM} 35(4), 1988.

\bibitem{Claude} Anthropic (Claude). A $2^{o(n)}$ time algorithm for 3-SAT. Prépublication,
 Hexagon 2610.00577v1, 26 octobre 2026.

\bibitem{DL} F. Dias, Y. Lima. Limits and constants of OpenAI's claimed ETH refutation.
 arXiv:2611.00517v1 [cs.CC], 31 octobre 2026.

\bibitem{DKSS13} Z. Dvir, S. Kopparty, S. Saraf, M. Sudan. Extensions to the method of
 multiplicities, with applications to Kakeya sets and mergers. \emph{SIAM J. Comput.} 42(6),
 2013.

\bibitem{Erd45} P. Erd\H{o}s. On a lemma of Littlewood and Offord. \emph{Bull. Amer. Math.
 Soc.} 51, 1945.

\bibitem{GLR82} I. Gohberg, P. Lancaster, L. Rodman. \emph{Matrix Polynomials}. Academic
 Press, 1982.

\bibitem{Hoe63} W. Hoeffding. Probability inequalities for sums of bounded random variables.
 \emph{J. Amer. Statist. Assoc.} 58(301), 1963.

\bibitem{IP01} R. Impagliazzo, R. Paturi. On the complexity of $k$-SAT. \emph{J. Comput.
 System Sci.} 62(2), 2001.

\bibitem{IPZ01} R. Impagliazzo, R. Paturi, F. Zane. Which problems have strongly exponential
 complexity? \emph{J. Comput. System Sci.} 63(4), 2001.

\bibitem{Kol58} A. Kolmogorov. Sur les propri\'et\'es des fonctions de concentrations de
 M.~P.~L\'evy. \emph{Ann. Inst. H. Poincar\'e} 16, 1958.

\bibitem{LT79} R. J. Lipton, R. E. Tarjan. A separator theorem for planar graphs.
 \emph{SIAM J. Appl. Math.} 36(2), 1979.

\bibitem{Lit30} J. E. Littlewood. On bounded bilinear forms in an infinite number of
 variables. \emph{Quart. J. Math.} 1, 1930.

\bibitem{LO43} J. E. Littlewood, A. C. Offord. On the number of real roots of a random
 algebraic equation (III). \emph{Rec. Math. [Mat. Sbornik] N.S.} 12(54), 1943.

\bibitem{Mathlib} The mathlib Community. The Lean mathematical library. \emph{CPP} 2020.

\bibitem{Lean4} L. de Moura, S. Ullrich. The Lean 4 theorem prover and programming language.
 \emph{CADE} 2021.

\bibitem{ODonnell} R. O'Donnell. Some exposition of ``The Exponential Time Hypothesis Is
 False''. Manuscrit, 2026.
 \url{https://www.cs.cmu.edu/~odonnell/some_eth_false_exposition.pdf}

\bibitem{OpenAI} OpenAI. The exponential time hypothesis is false. Prépublication, 17 octobre 2026,
 révisée le 25 octobre 2026.

\bibitem{RS86} N. Robertson, P. D. Seymour. Graph minors. II. Algorithmic aspects of
 tree-width. \emph{J. Algorithms} 7(3), 1986.

\bibitem{Rog61} B. A. Rogozin. On the increase of dispersion of sums of independent random
 variables. \emph{Theory Probab. Appl.} 6(1), 1961.

\bibitem{Tov84} C. A. Tovey. A simplified NP-complete satisfiability problem.
 \emph{Discrete Appl. Math.} 8, 1984.

\end{thebibliography}

\end{document}
